\documentclass[8pt]{beamer}
\usetheme{Madrid}
\usecolortheme{default}
\usepackage{amsmath, amssymb, graphicx, array, multirow, booktabs, xcolor, tikz}
\usetikzlibrary{shapes,arrows,positioning,fit,backgrounds}

% Compact font settings
\setbeamerfont{title}{size=\Large}
\setbeamerfont{subtitle}{size=\small}
\setbeamerfont{author}{size=\scriptsize}
\setbeamerfont{date}{size=\scriptsize}
\setbeamerfont{frametitle}{size=\large}
\setbeamerfont{framesubtitle}{size=\normalsize}
\setbeamerfont{enumerate item}{size=\footnotesize}
\setbeamerfont{itemize item}{size=\footnotesize}
\setbeamerfont{block title}{size=\footnotesize}

% Compact spacing
\setlength{\itemsep}{0.08ex}
\setlength{\parskip}{0.08ex}
\setlength{\leftmargini}{0.3cm}
\setbeamersize{text margin left=3mm, text margin right=3mm}
\addtobeamertemplate{frame begin}{\vspace{-0.4cm}}{}

% Colors
\definecolor{myblue}{RGB}{0,0,255}
\definecolor{myred}{RGB}{255,0,0}
\definecolor{mygreen}{RGB}{0,128,0}
\definecolor{myorange}{RGB}{255,165,0}
\definecolor{mypurple}{RGB}{128,0,128}
\definecolor{mygray}{RGB}{128,128,128}
\definecolor{mygold}{RGB}{255,215,0}

\title{Module 1: AI and Risk - Introduction and Overview}
\subtitle{100 Tutorial Topics: Classical AI, Neural Networks, Machine Learning Types, AI Risks}
\author{GARP RAI Exam Preparation}
\date{}

\begin{document}
	
	\begin{frame}
		\titlepage
	\end{frame}
	
	% ============================================================
	% SECTION 1: CLASSICAL AI
	% ============================================================
	
	\begin{frame}{T1: Introduction to Classical AI}
		\fontsize{6.5}{7.5}\selectfont
		\textbf{TOPIC: What is Classical AI?}
		
		\vspace{0.1cm}
		\textbf{DEFINITION:}
		\begin{itemize}
			\item The original paradigm of AI research
			\item Based on explicit, precise processing
			\item Typified by symbolic logic and formal rule-governed systems
			\item Also known as "Good Old-Fashioned AI" (GOFAI)
		\end{itemize}
		
		\vspace{0.1cm}
		\textbf{HISTORICAL ROOTS:}
		\begin{itemize}
			\item Charles Babbage (early 19th century): Difference engine and analytical engine
			\item Alan Turing (1936): Theoretical "Turing machine"
			\item Turing's 1950 paper "Computing Machinery and Intelligence"
			\item The "imitation game" (Turing Test)
		\end{itemize}
		
		\vspace{0.1cm}
		\textbf{KEY CHARACTERISTICS:}
		\begin{itemize}
			\item Intelligence as explicit rule-following
			\item Symbolic representation
			\item Logical deduction
			\item Search through possibilities
		\end{itemize}
		
		\textcolor{myred}{\textbf{EXAM TRAP:}} Classical AI = \textcolor{myblue}{explicit, rule-based processing}!
	\end{frame}
	
	\begin{frame}{T2: Turing Test}
		\fontsize{6.5}{7.5}\selectfont
		\textbf{TOPIC: The Turing Test (Imitation Game)}
		
		\vspace{0.1cm}
		\textbf{DEFINITION:}
		\begin{itemize}
			\item Proposed by Alan Turing in his 1950 paper
			\item Originally called the "imitation game"
			\item A test of machine intelligence
		\end{itemize}
		
		\vspace{0.1cm}
		\textbf{THE TEST:}
		\begin{itemize}
			\item A computer proves itself intelligent if it can generate textual conversation
			\item Quality indistinguishable from an intelligent human
			\item Ranging over any topic chosen by the human "interrogator"
		\end{itemize}
		
		\vspace{0.1cm}
		\textbf{SIGNIFICANCE:}
		\begin{itemize}
			\item Thought experiment demonstrating possibility of AI
			\item Highly controversial philosophically
			\item Remains influential
			\item Particularly salient with modern chatbots (ChatGPT)
		\end{itemize}
		
		\vspace{0.1cm}
		\textbf{LIMITATIONS:}
		\begin{itemize}
			\item Focuses on conversational ability, not general intelligence
			\item Can be gamed by sophisticated mimicry
			\item Doesn't require understanding, only imitation
		\end{itemize}
		
		\textcolor{myred}{\textbf{EXAM TRAP:}} Turing Test = \textcolor{myblue}{indistinguishable conversation}, not true understanding!
	\end{frame}
	
	\begin{frame}{T3: Specific versus General AI}
		\fontsize{6.5}{7.5}\selectfont
		\textbf{TOPIC: Specific versus General AI}
		
		\vspace{0.1cm}
		\textbf{SPECIFIC (NARROW) AI:}
		\begin{itemize}
			\item Designed to be intelligent at solving specific problems
			\item Outperforms humans in specific domains
			\item Examples: calculating tables, generating statistics, playing chess
			\item Most practical AI applications
		\end{itemize}
		
		\vspace{0.1cm}
		\textbf{GENERAL AI (AGI):}
		\begin{itemize}
			\item Comparable to human general intelligence
			\item Can handle any intellectual task
			\item Source of existential risk concerns
			\item Currently theoretical
		\end{itemize}
		
		\vspace{0.1cm}
		\textbf{KEY INSIGHT:}
		\begin{itemize}
			\item A system can be "intelligent" in specific tasks without general competence
			\item Turing's test was a conceptual existence proof
			\item Practical AI focuses on specific problem-solving
		\end{itemize}
		
		\textcolor{myred}{\textbf{EXAM TRAP:}} Most AI is \textcolor{myblue}{narrow/specific}, not general!
	\end{frame}
	
	\begin{frame}{T4: Dartmouth Conference}
		\fontsize{6.5}{7.5}\selectfont
		\textbf{TOPIC: The Dartmouth Conference (1956)}
		
		\vspace{0.1cm}
		\textbf{KEY FACTS:}
		\begin{itemize}
			\item Term "artificial intelligence" coined by John McCarthy
			\item Organized by McCarthy, Minsky, Rochester, and Shannon
			\item Two-month summer conference at Dartmouth College
			\item Seminal event in AI history
		\end{itemize}
		
		\vspace{0.1cm}
		\textbf{CORE CONJECTURE:}
		\begin{itemize}
			\item "Every aspect of learning or any other feature of intelligence can in principle be so precisely described that a machine can be made to simulate it"
			\item Intelligence as explicit, precise processing
			\item Symbolic logic as paradigm
		\end{itemize}
		
		\vspace{0.1cm}
		\textbf{SIGNIFICANCE:}
		\begin{itemize}
			\item Established AI as a research field
			\item Set the agenda for decades
			\item Echoed Turing's notion of universality
		\end{itemize}
		
		\textcolor{myred}{\textbf{EXAM TRAP:}} Dartmouth Conference (1956) = \textcolor{myblue}{birth of AI as a field}!
	\end{frame}
	
	\begin{frame}{T5: Logic Theorist}
		\fontsize{6.5}{7.5}\selectfont
		\textbf{TOPIC: Logic Theorist Program}
		
		\vspace{0.1cm}
		\textbf{KEY FACTS:}
		\begin{itemize}
			\item Developed by Newell, Simon, and Shaw (1955)
			\item Designed to generate proofs in propositional logic
			\item Presented at Dartmouth conference
			\item Confirmed most theorems in Principia Mathematica
		\end{itemize}
		
		\vspace{0.1cm}
		\textbf{ACHIEVEMENTS:}
		\begin{itemize}
			\item Found new and more elegant proof of Theorem 2.85
			\item Demonstrated machine intelligence in logic
			\item High-profile application of AI
		\end{itemize}
		
		\vspace{0.1cm}
		\textbf{SIGNIFICANCE:}
		\begin{itemize}
			\item Early success of classical AI
			\item Validated logicist thesis (mathematics from logic)
			\item Demonstrated rule-based reasoning
		\end{itemize}
		
		\textcolor{myred}{\textbf{EXAM TRAP:}} Logic Theorist = \textcolor{myblue}{automated theorem proving}!
	\end{frame}
	
	\begin{frame}{T6: Games as AI Testbeds}
		\fontsize{6.5}{7.5}\selectfont
		\textbf{TOPIC: Why Games Were Attractive AI Testbeds}
		
		\vspace{0.1cm}
		\textbf{THREE KEY REASONS:}
		\begin{enumerate}
			\item \textcolor{myblue}{Relatively Simple:}
			\begin{itemize}
				\item Compared to complexities of normal life
				\item Could be tackled with limited computing resources
				\item No need to record physical environment details
			\end{itemize}
			\item \textcolor{myblue}{Straightforwardly Rule-Governed:}
			\begin{itemize}
				\item Clear and explicit regulations
				\item Setup of the game
				\item How pieces are permitted to move
				\item Effects of moves on board position
				\item Criteria for winning, losing, drawing
			\end{itemize}
			\item \textcolor{myblue}{Competitive:}
			\begin{itemize}
				\item Players can be rated by relative skill
				\item Natural assessment of progress
				\item Straightforward to judge improvement
			\end{itemize}
		\end{enumerate}
		
		\textcolor{myred}{\textbf{EXAM TRAP:}} Games provided \textcolor{myblue}{controlled environments} for AI development!
	\end{frame}
	
	\begin{frame}{T7: Reinforcement Learning - Nim Example}
		\fontsize{6.5}{7.5}\selectfont
		\textbf{TOPIC: Simple Reinforcement Learning - Nim}
		
		\vspace{0.1cm}
		\textbf{THE GAME OF NIM:}
		\begin{itemize}
			\item 20 coins on a table
			\item Players alternate removing 1, 2, or 3 coins
			\item Winner removes the last coin(s)
			\item Winning strategy: reduce to multiple of 4
		\end{itemize}
		
		\vspace{0.1cm}
		\textbf{LEARNING MECHANISM:}
		\begin{itemize}
			\item Assign weights to each move in each position
			\item Initially all weights equal (e.g., 5)
			\item Choose moves randomly based on weights
			\item After each game, adjust weights based on outcome
		\end{itemize}
		
		\vspace{0.1cm}
		\textbf{REINFORCEMENT RULES:}
		\begin{itemize}
			\item \textcolor{myblue}{Winner:} Chosen moves that led to victory are positively reinforced (weight +1)
			\item \textcolor{myblue}{Winner:} Moves not chosen are negatively reinforced (weight -1)
			\item \textcolor{myblue}{Loser:} Chosen moves that led to loss are negatively reinforced
			\item \textcolor{myblue}{Loser:} Moves not chosen are positively reinforced
		\end{itemize}
		
		\vspace{0.1cm}
		\textbf{RESULT:}
		\begin{itemize}
			\item After ~100 games, discovers winning strategy
			\item Learns to reduce coins to multiples of 4
			\item No explicit "lookahead" or planning
		\end{itemize}
		
		\textcolor{myred}{\textbf{EXAM TRAP:}} Reinforcement learning = \textcolor{myblue}{learning from trial and error}!
	\end{frame}
	
	\begin{frame}{T8: Lookahead}
		\fontsize{6.5}{7.5}\selectfont
		\textbf{TOPIC: Lookahead in AI}
		
		\vspace{0.1cm}
		\textbf{DEFINITION:}
		\begin{itemize}
			\item Forming a multistep plan to achieve a goal
			\item Exploring consequences of possible actions
			\item Essential when many possible situations exist
		\end{itemize}
		
		\vspace{0.1cm}
		\textbf{WHEN LOOKAHEAD IS ESSENTIAL:}
		\begin{itemize}
			\item Very large number of possible situations overall
			\item Relatively constrained and predictable options in particular situations
			\item Need to form multistep plans
		\end{itemize}
		
		\vspace{0.1cm}
		\textbf{EXAMPLES:}
		\begin{itemize}
			\item \textcolor{myblue}{Sudoku:} "If I put 4 here, then this cell must be 7..."
			\item \textcolor{myblue}{Tower of Hanoi:} Strategic planning to move disks
			\item \textcolor{myblue}{Chess:} Considering future moves and responses
		\end{itemize}
		
		\vspace{0.1cm}
		\textbf{APPROACHES:}
		\begin{itemize}
			\item Tactical lookahead: Narrow down plausible sequences
			\item Strategic planning: Identify useful shorter-term outcomes
		\end{itemize}
		
		\textcolor{myred}{\textbf{EXAM TRAP:}} Lookahead = \textcolor{myblue}{exploring future consequences} of actions!
	\end{frame}
	
	\begin{frame}{T9: Search Techniques}
		\fontsize{6.5}{7.5}\selectfont
		\textbf{TOPIC: Search in Classical AI}
		
		\vspace{0.1cm}
		\textbf{DEFINITION:}
		\begin{itemize}
			\item Generating sets of possible outcomes
			\item Searching through them for best achievement of aims
			\item Fundamental technique in AI
		\end{itemize}
		
		\vspace{0.1cm}
		\textbf{BINARY SEARCH:}
		\begin{itemize}
			\item Progressively divide search space in two
			\item Used in number-guessing program
			\item Finds target within at most 10 guesses for 1,000 possibilities
			\item Efficient: $2^{10} = 1,024$
		\end{itemize}
		
		\vspace{0.1cm}
		\textbf{ROUTE-FINDING SEARCH:}
		\begin{itemize}
			\item \textcolor{myblue}{Breadth-First Search:} Track all possible routes, build stage by stage
			\item \textcolor{myblue}{Depth-First Search:} Follow each route to end before comparing next
			\item \textcolor{myblue}{Heuristic Search:} Use additional information to guide exploration
		\end{itemize}
		
		\vspace{0.1cm}
		\textbf{A* ALGORITHM:}
		\begin{itemize}
			\item Best-known heuristic search technique
			\item Tracks distance traveled so far
			\item Estimates distance still to go (optimistic)
			\item Orders locations by sum of these values
		\end{itemize}
		
		\textcolor{myred}{\textbf{EXAM TRAP:}} Search = \textcolor{myblue}{systematic exploration of possibilities}!
	\end{frame}
	
	\begin{frame}{T10: Combinatorial Explosion}
		\fontsize{6.5}{7.5}\selectfont
		\textbf{TOPIC: Combinatorial Explosion}
		
		\vspace{0.1cm}
		\textbf{DEFINITION:}
		\begin{itemize}
			\item Search space grows enormously as problem size increases
			\item Relevant numbers of possible choices get multiplied together
			\item Certainty of solution quickly becomes unfeasible
		\end{itemize}
		
		\vspace{0.1cm}
		\textbf{EXAMPLES:}
		\begin{itemize}
			\item \textcolor{myblue}{Travelling Salesman Problem:} Finding optimal route around cities
			\item \textcolor{myblue}{Chess:} Average branching factor ~30
			\item \textcolor{myblue}{Go:} Average branching factor ~250
			\item \textcolor{myblue}{Checkers:} Average branching factor ~10
		\end{itemize}
		
		\vspace{0.1cm}
		\textbf{CONSEQUENCES:}
		\begin{itemize}
			\item Exhaustive search impractical for complex problems
			\item Earth would be swallowed by Sun before first move
			\item Need for heuristics and approximations
			\item Adversarial search adds further complications
		\end{itemize}
		
		\textcolor{myred}{\textbf{EXAM TRAP:}} Combinatorial explosion = \textcolor{myblue}{exponential growth in possibilities}!
	\end{frame}
	
	\begin{frame}{T11: Recursion}
		\fontsize{6.5}{7.5}\selectfont
		\textbf{TOPIC: Recursion in AI}
		
		\vspace{0.1cm}
		\textbf{DEFINITION:}
		\begin{itemize}
			\item Powerful technique in classical AI
			\item Complex problem solved by reducing to simpler instances of same problem
			\item Huge importance in computer science generally
		\end{itemize}
		
		\vspace{0.1cm}
		\textbf{TOWER OF HANOI EXAMPLE:}
		\begin{itemize}
			\item Move pile of disks from pillar A to pillar B
			\item Break down into three subtasks:
			\begin{enumerate}
				\item Move pile of n-1 disks from A to C
				\item Move disk n from A to B
				\item Move pile of n-1 disks from C to B
			\end{enumerate}
			\item Each subtask reduces to simpler instance
			\item Eventually reduces to single disk movements
		\end{itemize}
		
		\vspace{0.1cm}
		\textbf{MATHEMATICAL RESULT:}
		\begin{itemize}
			\item Moving n disks takes $2^n - 1$ steps
			\item Moving 9 disks takes 511 steps
			\item Proved by mathematical induction
		\end{itemize}
		
		\textcolor{myred}{\textbf{EXAM TRAP:}} Recursion = \textcolor{myblue}{solving problems by reducing to simpler versions}!
	\end{frame}
	
	\begin{frame}{T12: Recursive Adversarial Tree Search}
		\fontsize{6.5}{7.5}\selectfont
		\textbf{TOPIC: Recursive Adversarial Tree Search - Noughts and Crosses}
		
		\vspace{0.1cm}
		\textbf{GOAL:}
		\begin{itemize}
			\item Design program to play perfect game of Noughts-and-Crosses
			\item Never lose
			\item Always find winning sequence if one exists
		\end{itemize}
		
		\vspace{0.1cm}
		\textbf{POSITION VALUES:}
		\begin{itemize}
			\item \textcolor{myblue}{+1:} Win for player whose turn it is
			\item \textcolor{myblue}{-1:} Loss for player whose turn it is
			\item \textcolor{myblue}{0:} Draw
		\end{itemize}
		
		\vspace{0.1cm}
		\textbf{ALGORITHM:}
		\begin{enumerate}
			\item If X has lost (opponent has line of three), value = -1
			\item If position is full, value = 0 (draw)
			\item Otherwise, construct all positions with Y to move
			\item Evaluate each from Y's point of view using same algorithm
			\item Assign X-value = inverse of lowest Y-value
		\end{enumerate}
		
		\vspace{0.1cm}
		\textbf{MINIMAX SEARCH:}
		\begin{itemize}
			\item Presume opponent chooses reply maximizing their benefit
			\item Choose move minimizing that maximum benefit
			\item Applied recursively at every level
		\end{itemize}
		
		\textcolor{myred}{\textbf{EXAM TRAP:}} Adversarial search = \textcolor{myblue}{minimax principle}!
	\end{frame}
	
	\begin{frame}{T13: Heuristics and Refinement}
		\fontsize{6.5}{7.5}\selectfont
		\textbf{TOPIC: Heuristics and Refinement by Reinforcement Learning}
		
		\vspace{0.1cm}
		\textbf{HEURISTICS DEFINITION:}
		\begin{itemize}
			\item Calculable criteria or rules of thumb
			\item Guide choice of position to aim for
			\item Cannot guarantee optimal solution
		\end{itemize}
		
		\vspace{0.1cm}
		\textbf{CHECKERS EXAMPLE - CRITERIA:}
		\begin{enumerate}
			\item Material balance (my pieces vs. opponent's)
			\item Relative mobility of pieces
			\item How far advanced pieces are
			\item Relative number of "kings"
			\item Additional criteria (Samuel devised 38)
		\end{enumerate}
		
		\vspace{0.1cm}
		\textbf{WEIGHTED EVALUATION:}
		\[
		v = w_1 c_1 + w_2 c_2 + w_3 c_3 + w_4 c_4 + w_5 c_5 + w_6 c_6
		\]
		
		\vspace{0.1cm}
		\textbf{REFINEMENT BY REINFORCEMENT LEARNING:}
		\begin{itemize}
			\item Initial weights equal (e.g., 1,000)
			\item Random variation to create competing strategy
			\item Play games between current and competing
			\item Replace current with competing if better
			\item Repeat iteratively
		\end{itemize}
		
		\textcolor{myred}{\textbf{EXAM TRAP:}} Heuristics guide search; \textcolor{myblue}{weights refined by reinforcement learning}!
	\end{frame}
	
	\begin{frame}{T14: Limits of Classical AI}
		\fontsize{6.5}{7.5}\selectfont
		\textbf{TOPIC: Limitations of Classical AI}
		
		\vspace{0.1cm}
		\textbf{KEY DIFFICULTIES:}
		\begin{enumerate}
			\item \textcolor{myblue}{Representation Problem:}
			\begin{itemize}
				\item How to represent characteristics and relationships
				\item Especially rules governing behavior
				\item Experts find it hard to elucidate implicit assumptions
			\end{itemize}
			\item \textcolor{myblue}{Frame Problem:}
			\begin{itemize}
				\item Tracking what changes when actions performed
				\item Also tracking what stays the same
				\item Adds to combinatorial explosion
			\end{itemize}
			\item \textcolor{myblue}{Uncertainty:}
			\begin{itemize}
				\item Ambiguous or uncertain data
				\item Complex probability calculations
				\item Unlike perfect-information games
			\end{itemize}
			\item \textcolor{myblue}{Real-Time Data:}
			\begin{itemize}
				\item Error-prone acquisition
				\item Computer vision challenges
				\item Object recognition difficulties
			\end{itemize}
		\end{enumerate}
		
		\textcolor{myred}{\textbf{EXAM TRAP:}} Classical AI is \textcolor{myblue}{brittle - fails beyond narrow boundaries}!
	\end{frame}
	
	\begin{frame}{T15: Expert Systems}
		\fontsize{6.5}{7.5}\selectfont
		\textbf{TOPIC: Expert Systems}
		
		\vspace{0.1cm}
		\textbf{DEFINITION:}
		\begin{itemize}
			\item Designed to capture knowledge within specific domains
			\item Often use logic programming
			\item Response to limitations of general-purpose AI
			\item Especially from 1980s
		\end{itemize}
		
		\vspace{0.1cm}
		\textbf{ARCHITECTURE:}
		\begin{itemize}
			\item \textcolor{myblue}{Knowledge Base:} Topic-specific facts and rules
			\item \textcolor{myblue}{Inference Engine:} General-purpose reasoning mechanism
			\item Knowledge engineers extend to new domains
		\end{itemize}
		
		\vspace{0.1cm}
		\textbf{EXAMPLE - JAPANESE FIFTH GENERATION:}
		\begin{itemize}
			\item Massive government investment
			\item Harness PROLOG as inferential medium
			\item Aimed for general-purpose reasoning
		\end{itemize}
		
		\vspace{0.1cm}
		\textbf{LIMITATIONS:}
		\begin{itemize}
			\item Functionality not easily generalized
			\item Hit same old problems beyond specific domains
			\item Brittle - fail unexpectedly beyond boundaries
		\end{itemize}
		
		\textcolor{myred}{\textbf{EXAM TRAP:}} Expert systems = \textcolor{myblue}{domain-specific knowledge capture}!
	\end{frame}
	
	% ============================================================
	% SECTION 2: NEURAL NETWORKS
	% ============================================================
	
	\begin{frame}{T16: Early Neural Network History}
		\fontsize{6.5}{7.5}\selectfont
		\textbf{TOPIC: Early History of Neural Networks}
		
		\vspace{0.1cm}
		\textbf{McCULLOCH AND PITTS (1943):}
		\begin{itemize}
			\item "A logical calculus of the ideas immanent in nervous activity"
			\item Described activity of neurons in brain
			\item Proposed networks analyzed in terms of propositional logic
			\item Neural nets equivalent to Turing machines
		\end{itemize}
		
		\vspace{0.1cm}
		\textbf{DONALD HEBB (1949):}
		\begin{itemize}
			\item Book: "The Organisation of Behaviour"
			\item Inspired by biology rather than mathematics
			\item Suggested associative learning in neural networks
			\item "Neurophysiological postulate" about connection strength
		\end{itemize}
		
		\vspace{0.1cm}
		\textbf{FRANK ROSENBLATT (1958):}
		\begin{itemize}
			\item Developed the Perceptron model
			\item Rejected symbolic/algorithmic approach
			\item Formulated in terms of probability theory
			\item Inspired widespread interest in neural networks
		\end{itemize}
		
		\textcolor{myred}{\textbf{EXAM TRAP:}} Early neural networks = \textcolor{myblue}{biologically inspired learning}!
	\end{frame}
	
	\begin{frame}{T17: Artificial Neurons}
		\fontsize{6.5}{7.5}\selectfont
		\textbf{TOPIC: Artificial Neurons}
		
		\vspace{0.1cm}
		\textbf{STRUCTURE:}
		\begin{itemize}
			\item Takes numerical inputs ($i_1, i_2, i_3$, etc.)
			\item Each input multiplied by weight ($w_1, w_2, w_3$, etc.)
			\item Weights added together (plus bias term $w_0$)
			\item Net input fed into activation function $f$
			\item Output value $v$ determined
		\end{itemize}
		
		\vspace{0.1cm}
		\textbf{FORMULA:}
		\[
		v = f(w_0 + w_1 i_1 + w_2 i_2 + w_3 i_3 + ...)
		\]
		
		\vspace{0.1cm}
		\textbf{ACTIVATION FUNCTION:}
		\begin{itemize}
			\item Chosen to be non-linear
			\item Often approximates step function
			\item Triggered when net input > threshold
			\item Non-linearity crucial for learning non-linear behavior
		\end{itemize}
		
		\vspace{0.1cm}
		\textbf{COMMON ACTIVATION FUNCTIONS:}
		\begin{itemize}
			\item Sigmoid
			\item Hyperbolic tangent (tanh)
			\item Rectified Linear Unit (ReLU)
		\end{itemize}
		
		\textcolor{myred}{\textbf{EXAM TRAP:}} Artificial neuron = \textcolor{myblue}{weighted sum + non-linear activation}!
	\end{frame}
	
	\begin{frame}{T18: Deep Neural Networks}
		\fontsize{6.5}{7.5}\selectfont
		\textbf{TOPIC: Deep Neural Networks}
		
		\vspace{0.1cm}
		\textbf{STRUCTURE:}
		\begin{itemize}
			\item Input layer (I)
			\item Multiple hidden layers (a, b, c, etc.)
			\item Output layer (O)
			\item Each neuron in one layer connected to all in next
		\end{itemize}
		
		\vspace{0.1cm}
		\textbf{HOW IT WORKS:}
		\begin{itemize}
			\item Input layer: specific numeric activation levels represent input data
			\item Example: pixel color values in an image
			\item Hidden layers: activation depends on previous layer
			\item Weights vary, influencing propagation of activity
			\item Learning = changing weights
		\end{itemize}
		
		\vspace{0.1cm}
		\textbf{EXAMPLE - DIGIT RECOGNITION:}
		\begin{itemize}
			\item Input: pixel values of digit image
			\item Output: 10 neurons (one per digit 0-9)
			\item Correct digit: activation close to 1
			\item Other digits: activation close to 0
		\end{itemize}
		
		\textcolor{myred}{\textbf{EXAM TRAP:}} Deep networks = \textcolor{myblue}{multiple hidden layers} between input and output!
	\end{frame}
	
	\begin{frame}{T19: Minsky and Papert Critique}
		\fontsize{6.5}{7.5}\selectfont
		\textbf{TOPIC: Minsky and Papert's Critique}
		
		\vspace{0.1cm}
		\textbf{THE BOOK (1969):}
		\begin{itemize}
			\item "Perceptrons" by Marvin Minsky and Seymour Papert
			\item Fierce critique of single-layer Perceptrons
			\item Severely undermined enthusiasm for neural networks
		\end{itemize}
		
		\vspace{0.1cm}
		\textbf{KEY ARGUMENT:}
		\begin{itemize}
			\item Perceptrons incapable of learning simple logical functions
			\item Notably anything depending on "exclusive or" (XOR)
			\item Appeared to wreck prospect of intelligent information processing
		\end{itemize}
		
		\vspace{0.1cm}
		\textbf{LATER REALIZATION:}
		\begin{itemize}
			\item Critique applied only to single-layer networks
			\item More complex functions learnable with multiple layers
			\item But adding layers complicated learning process
			\item Route largely ignored until backpropagation in 1980s
		\end{itemize}
		
		\textcolor{myred}{\textbf{EXAM TRAP:}} Minsky-Papert critique applied \textcolor{myblue}{only to single-layer networks}!
	\end{frame}
	
	\begin{frame}{T20: Backpropagation}
		\fontsize{6.5}{7.5}\selectfont
		\textbf{TOPIC: Backpropagation Learning Algorithm}
		
		\vspace{0.1cm}
		\textbf{PURPOSE:}
		\begin{itemize}
			\item Learning algorithm for multi-layer neural networks
			\item Rose to prominence in 1980s
			\item Enabled training of deep networks
		\end{itemize}
		
		\vspace{0.1cm}
		\textbf{HOW IT WORKS:}
		\begin{enumerate}
			\item Input data (e.g., image pixels)
			\item Compare output activation to desired
			\item Calculate how small change in output layer weights would improve match
			\item Work back through layers (backpropagation)
			\item Calculate for all weights in network
			\item Adjust weights by small amount in direction of improvement
			\item Repeat with other images (thousands/millions of times)
		\end{enumerate}
		
		\vspace{0.1cm}
		\textbf{TECHNICAL BASIS:}
		\begin{itemize}
			\item Chain rule of differentiation
			\item Partial derivatives of loss function
			\item Iterative weight adjustment
		\end{itemize}
		
		\textcolor{myred}{\textbf{EXAM TRAP:}} Backpropagation = \textcolor{myblue}{working backward through layers} to adjust weights!
	\end{frame}
	
	\begin{frame}{T21: Connectionism}
		\fontsize{6.5}{7.5}\selectfont
		\textbf{TOPIC: Connectionism}
		
		\vspace{0.1cm}
		\textbf{DEFINITION:}
		\begin{itemize}
			\item Approach emphasizing neural networks
			\item Displaced pessimism of Minsky and Papert
			\item Demonstrated learning of complex functions
		\end{itemize}
		
		\vspace{0.1cm}
		\textbf{KEY WORK (1986):}
		\begin{itemize}
			\item "Parallel Distributed Processing" by Rumelhart, McClelland, and PDP Research Group
			\item Two volumes
			\item Showed how layers of simple interacting neurons achieve learning
		\end{itemize}
		
		\vspace{0.1cm}
		\textbf{FOUR KEY ATTRACTIONS:}
		\begin{enumerate}
			\item \textcolor{myblue}{Biologically Inspired:} Potential to shed light on how we think
			\item \textcolor{myblue}{Natural Learning:} Through association and feedback
			\item \textcolor{myblue}{General:} Applied to huge variety of tasks without special programming
			\item \textcolor{myblue}{Distributed:} Robust to noisy data, generalizes well, graceful degradation
		\end{enumerate}
		
		\vspace{0.1cm}
		\textbf{LATER DECLINE:}
		\begin{itemize}
			\item Enthusiasm steeply declined in 1990s
			\item Failed to deliver expected practical breakthroughs
			\item Day would come with greater computational power
		\end{itemize}
		
		\textcolor{myred}{\textbf{EXAM TRAP:}} Connectionism = \textcolor{myblue}{neural network approach} to cognition!
	\end{frame}
	
	\begin{frame}{T22: LeNet-5}
		\fontsize{6.5}{7.5}\selectfont
		\textbf{TOPIC: LeNet-5 - Convolutional Neural Network}
		
		\vspace{0.1cm}
		\textbf{KEY FACTS:}
		\begin{itemize}
			\item Developed by Yann LeCun and colleagues (1998)
			\item Designed for recognition of digits
			\item 8 layers
			\item Convolutional neural network
		\end{itemize}
		
		\vspace{0.1cm}
		\textbf{CONVOLUTIONAL LAYERS:}
		\begin{itemize}
			\item Apply local filter (kernel) repeatedly across grid
			\item Filter = small matrix of weights
			\item Weights multiplied by pixel activation values
			\item Sum of products = activation in convolutional layer
			\item Weights learned like other network weights
		\end{itemize}
		
		\vspace{0.1cm}
		\textbf{ARCHITECTURE:}
		\begin{itemize}
			\item Input: 32 x 32 pixel image
			\item Two sequences: convolution + pooling
			\item Pooling reduces dimensions (averaging)
			\item Final layers: 120, 84, 10 neurons
			\item Total: 9,118 neurons
		\end{itemize}
		
		\textcolor{myred}{\textbf{EXAM TRAP:}} Convolutional layers = \textcolor{myblue}{efficient local feature identification}!
	\end{frame}
	
	\begin{frame}{T23: AlexNet Breakthrough}
		\fontsize{6.5}{7.5}\selectfont
		\textbf{TOPIC: AlexNet - 2012 ImageNet Breakthrough}
		
		\vspace{0.1cm}
		\textbf{KEY FACTS:}
		\begin{itemize}
			\item Developed by Krizhevsky, Sutskever, and Hinton (2012)
			\item Won ImageNet Large Scale Visual Recognition Challenge
			\item Error rate: 15.3\% (vs. 25.8\% in 2011)
			\item Nearly 100 times bigger than LeNet-5
			\item Almost 900,000 neurons
		\end{itemize}
		
		\vspace{0.1cm}
		\textbf{ENABLING TECHNOLOGY:}
		\begin{itemize}
			\item Graphics processing units (GPUs)
			\item Developed for animated computer games
			\item Process image pixels in parallel
			\item Recent versions: up to 10,000 cores
			\item Perfect for huge matrix operations in neural networks
		\end{itemize}
		
		\vspace{0.1cm}
		\textbf{SIGNIFICANCE:}
		\begin{itemize}
			\item Startling vindication of deep networks
			\item Demonstrated amazing power for image recognition
			\item Sparked renewed interest in deep learning
		\end{itemize}
		
		\textcolor{myred}{\textbf{EXAM TRAP:}} AlexNet (2012) = \textcolor{myblue}{breakthrough moment} for deep learning!
	\end{frame}
	
	\begin{frame}{T24: Atari Games and Deep Reinforcement Learning}
		\fontsize{6.5}{7.5}\selectfont
		\textbf{TOPIC: DeepMind's Atari Breakthrough}
		
		\vspace{0.1cm}
		\textbf{KEY FACTS (2013):}
		\begin{itemize}
			\item DeepMind (London company, founded 2010)
			\item Deep convolutional network
			\item Learned to play vintage Atari games
			\item Including Breakout and Pong
			\item Performed vastly better than expert human
		\end{itemize}
		
		\vspace{0.1cm}
		\textbf{REMARKABLE ASPECT:}
		\begin{itemize}
			\item No information about goal of games
			\item No information about significance of screen images
			\item No information about effect of user actions
			\item Learned entirely from trying actions
			\item Based on pixel information and game score
		\end{itemize}
		
		\vspace{0.1cm}
		\textbf{DEEP REINFORCEMENT LEARNING:}
		\begin{itemize}
			\item Based on score feedback
			\item System teaching itself from scratch
			\item Learned what to do entirely on its own
		\end{itemize}
		
		\textcolor{myred}{\textbf{EXAM TRAP:}} Atari program = \textcolor{myblue}{self-taught learning from scratch}!
	\end{frame}
	
	\begin{frame}{T25: AlphaGo and AlphaZero}
		\fontsize{6.5}{7.5}\selectfont
		\textbf{TOPIC: AlphaGo and AlphaZero}
		
		\vspace{0.1cm}
		\textbf{ALPHAGO (2014-2016):}
		\begin{itemize}
			\item DeepMind's deep-learning program
			\item Defeated European champion Fan Hui 5-0 (2014)
			\item Defeated World Champion Lee Sedol 4-1 (2015)
			\item Go considered too subtle for computer algorithms
		\end{itemize}
		
		\vspace{0.1cm}
		\textbf{ALPHAZERO (2017):}
		\begin{itemize}
			\item More generalized program
			\item Capable of teaching itself new games
			\item Including both Go and chess
			\item Given: starting position, moves, rules
			\item Built-in: recursive adversarial tree search
		\end{itemize}
		
		\vspace{0.1cm}
		\textbf{KEY DIFFERENCE:}
		\begin{itemize}
			\item Learned by competing against itself
			\item No input from human experts
			\item No game databases
			\item Self-trained in only a few hours
			\item Defeated world champion chess program Stockfish
		\end{itemize}
		
		\textcolor{myred}{\textbf{EXAM TRAP:}} AlphaZero = \textcolor{myblue}{self-teaching through self-play}!
	\end{frame}
	
	\begin{frame}{T26: Inscrutability of Deep Learning}
		\fontsize{6.5}{7.5}\selectfont
		\textbf{TOPIC: The Inscrutability of Deep Learning}
		
		\vspace{0.1cm}
		\textbf{DEFINITION:}
		\begin{itemize}
			\item Information stored implicitly within network weights
			\item Complex pattern of weights represents what network has learned
			\item Impossible for unaided human to interpret
			\item Behavior predictable only by experience
		\end{itemize}
		
		\vspace{0.1cm}
		\textbf{KEY CHARACTERISTICS:}
		\begin{itemize}
			\item Neither weights nor roles predetermined
			\item Intermediate layer representations learned
			\item No explanation of decisions available
			\item "Black box" nature
		\end{itemize}
		
		\vspace{0.1cm}
		\textbf{ANALOGY TO HUMAN COGNITION:}
		\begin{itemize}
			\item Closer to unconscious pattern-recognition
			\item Unlike explicit calculation or reasoning
			\item We also learn by examples without awareness
			\item Hard to articulate how we recognize things
		\end{itemize}
		
		\vspace{0.1cm}
		\textbf{IMPORTANT DISTINCTION:}
		\begin{itemize}
			\item Machines do NOT "understand" reflectively
			\item Completely non-sentient
			\item No awareness or conscious understanding
		\end{itemize}
		
		\textcolor{myred}{\textbf{EXAM TRAP:}} Inscrutability = \textcolor{myblue}{implicit, incomprehensible representations}!
	\end{frame}
	
	\begin{frame}{T27: ChatGPT and Large Language Models}
		\fontsize{6.5}{7.5}\selectfont
		\textbf{TOPIC: ChatGPT and Large Language Models}
		
		\vspace{0.1cm}
		\textbf{KEY FACTS (2022 onwards):}
		\begin{itemize}
			\item OpenAI's chatbot
			\item Developed using deep learning
			\item Trained on ~300 billion words from internet
			\item Potential to pass Turing Test in full generality
		\end{itemize}
		
		\vspace{0.1cm}
		\textbf{HOW IT WORKS:}
		\begin{itemize}
			\item Statistical analysis on textual data
			\item Records probability that word sequences will continue
			\item Predicts most likely words to follow
			\item Chooses words (with some randomness)
			\item Adds word to text, repeats
		\end{itemize}
		
		\vspace{0.1cm}
		\textbf{TRAINING PROCESS:}
		\begin{itemize}
			\item Automated learning: trillion parameters
			\item Supervised fine-tuning: human-crafted responses
			\item Reinforcement learning from human feedback
			\item Reward model rates responses
		\end{itemize}
		
		\vspace{0.1cm}
		\textbf{KEY INSIGHT:}
		\begin{itemize}
			\item Method based on imitating training data
			\item No internal model of domain
			\item "Stochastic parrot"
		\end{itemize}
		
		\textcolor{myred}{\textbf{EXAM TRAP:}} ChatGPT = \textcolor{myblue}{imitating textual patterns}, not understanding!
	\end{frame}
	
	\begin{frame}{T28: Limitations of ChatGPT}
		\fontsize{6.5}{7.5}\selectfont
		\textbf{TOPIC: Limitations of ChatGPT}
		
		\vspace{0.1cm}
		\textbf{CHESS EXAMPLE:}
		\begin{itemize}
			\item Common opening: responds with sensible moves
			\item Gives impression of understanding chessboard
			\item But: no internal model of board position
			\item No mechanisms of analysis or lookahead
			\item Subject to absurd errors:
			\begin{itemize}
				\item Overlooking obvious captures
				\item Allowing bishop to capture own piece
				\item Losing track of pawns
			\end{itemize}
		\end{itemize}
		
		\vspace{0.1cm}
		\textbf{WHY IT HAPPENS:}
		\begin{itemize}
			\item Stores textual patterns, not domain knowledge
			\item Apparent intelligence based on illusion
			\item No internal understanding of what's discussed
		\end{itemize}
		
		\vspace{0.1cm}
		\textbf{POTENTIAL FIX:}
		\begin{itemize}
			\item Link to dedicated chess engine
			\item Engine has model of board and analytical algorithms
			\item But general intelligence remains illusory
		\end{itemize}
		
		\textcolor{myred}{\textbf{EXAM TRAP:}} ChatGPT's "intelligence" = \textcolor{myblue}{illusion based on mimicry}!
	\end{frame}
	
	% ============================================================
	% SECTION 3: MACHINE LEARNING AND ITS RISKS
	% ============================================================
	
	\begin{frame}{T29: Four Types of Machine Learning}
		\fontsize{6.5}{7.5}\selectfont
		\textbf{TOPIC: Four Main Types of Machine Learning}
		
		\vspace{0.1cm}
		\textbf{1. REINFORCEMENT LEARNING:}
		\begin{itemize}
			\item Sequence of actions or decisions
			\item Changing environment
			\item Trial and error
			\item Positive/negative feedback
		\end{itemize}
		
		\vspace{0.1cm}
		\textbf{2. SUPERVISED LEARNING:}
		\begin{itemize}
			\item Labelled examples (training data)
			\item Categorize or predict
			\item Deep learning particularly risky
			\item Inscrutability and hidden biases
		\end{itemize}
		
		\vspace{0.1cm}
		\textbf{3. UNSUPERVISED LEARNING:}
		\begin{itemize}
			\item Identify patterns without labels
			\item Cluster analysis
			\item Dimensionality reduction
			\item Principal Component Analysis (PCA)
		\end{itemize}
		
		\vspace{0.1cm}
		\textbf{4. SEMI-SUPERVISED LEARNING:}
		\begin{itemize}
			\item Combination of supervised and unsupervised
			\item Only part of data labelled
			\item Unlabeled data for patterns
			\item Pseudo-labels
		\end{itemize}
		
		\textcolor{myred}{\textbf{EXAM TRAP:}} Four types = \textcolor{myblue}{reinforcement, supervised, unsupervised, semi-supervised}!
	\end{frame}
	
	\begin{frame}{T30: Reinforcement Learning Applications}
		\fontsize{6.5}{7.5}\selectfont
		\textbf{TOPIC: Reinforcement Learning Applications}
		
		\vspace{0.1cm}
		\textbf{DEFINITION:}
		\begin{itemize}
			\item Sequence of actions/decisions
			\item Changing environment
			\item Goal achievement
			\item Trial and error with feedback
		\end{itemize}
		
		\vspace{0.1cm}
		\textbf{KEY APPLICATIONS:}
		\begin{enumerate}
			\item \textcolor{myblue}{Robotics:}
			\begin{itemize}
				\item Teaching robots/autonomous vehicles
				\item Move and manipulate objects
				\item Interact with humans
				\item Learn from human feedback
			\end{itemize}
			\item \textcolor{myblue}{Healthcare:}
			\begin{itemize}
				\item Monitoring treatment effects
				\item Groups and individuals
			\end{itemize}
			\item \textcolor{myblue}{Finance:}
			\begin{itemize}
				\item Trading algorithms
			\end{itemize}
			\item \textcolor{myblue}{Optimization:}
			\begin{itemize}
				\item Any area where optimization by learning from experience
			\end{itemize}
		\end{enumerate}
		
		\vspace{0.1cm}
		\textbf{EXAMPLES FROM MODULE:}
		\begin{itemize}
			\item Nim game
			\item Atari games
			\item AlphaGo
			\item AlphaZero
		\end{itemize}
		
		\textcolor{myred}{\textbf{EXAM TRAP:}} Reinforcement learning = \textcolor{myblue}{learning from experience with feedback}!
	\end{frame}
	
	\begin{frame}{T31: Supervised Learning Applications}
		\fontsize{6.5}{7.5}\selectfont
		\textbf{TOPIC: Supervised Learning Applications}
		
		\vspace{0.1cm}
		\textbf{DEFINITION:}
		\begin{itemize}
			\item Learn from labelled examples
			\item Categorize or predict
			\item Training data with labels
		\end{itemize}
		
		\vspace{0.1cm}
		\textbf{CLASSICAL METHODS:}
		\begin{itemize}
			\item Linear regression
			\item Decision trees
			\item Support vector classification
		\end{itemize}
		
		\vspace{0.1cm}
		\textbf{DEEP LEARNING APPLICATIONS:}
		\begin{enumerate}
			\item \textcolor{myblue}{Image Classification:}
			\begin{itemize}
				\item Handwriting recognition
				\item Face recognition
				\item Road sign recognition
				\item Medical scans
				\item Visual scenes
			\end{itemize}
			\item \textcolor{myblue}{Speech Recognition}
			\item \textcolor{myblue}{Machine Translation}
			\item \textcolor{myblue}{Recommendation Systems}
		\end{enumerate}
		
		\vspace{0.1cm}
		\textbf{LEARNED PREDICTION (REGRESSION):}
		\begin{itemize}
			\item Numerical values rather than categories
			\item Financial analytics: stock prices, demand
			\item Asset valuation: house prices
			\item Credit scoring: probability of default
			\item Weather prediction
		\end{itemize}
		
		\textcolor{myred}{\textbf{EXAM TRAP:}} Supervised learning = \textcolor{myblue}{learning from labelled examples}!
	\end{frame}
	
	\begin{frame}{T32: Unsupervised Learning Applications}
		\fontsize{6.5}{7.5}\selectfont
		\textbf{TOPIC: Unsupervised Learning Applications}
		
		\vspace{0.1cm}
		\textbf{DEFINITION:}
		\begin{itemize}
			\item Identify patterns without labels
			\item No human input in form of labels
		\end{itemize}
		
		\vspace{0.1cm}
		\textbf{TWO PROMINENT FORMS:}
		\begin{enumerate}
			\item \textcolor{myblue}{Cluster Analysis:}
			\begin{itemize}
				\item Groups of similar objects identified
				\item Without training set
				\item Automatic identification
			\end{itemize}
			\item \textcolor{myblue}{Dimensionality Reduction:}
			\begin{itemize}
				\item Main dimensions of variation identified
				\item Range and relationships more easily grasped
				\item Comparative closeness
			\end{itemize}
		\end{enumerate}
		
		\vspace{0.1cm}
		\textbf{PRINCIPAL COMPONENT ANALYSIS (PCA):}
		\begin{itemize}
			\item Classical method
			\item Illustrates both cluster analysis and dimensionality reduction
			\item Application to village layout
			\item Application to text analysis
		\end{itemize}
		
		\vspace{0.1cm}
		\textbf{KEY CHARACTERISTIC:}
		\begin{itemize}
			\item Less inscrutable than deep learning
			\item Can understand what dimensions represent
			\item Explore basis relatively straightforwardly
		\end{itemize}
		
		\textcolor{myred}{\textbf{EXAM TRAP:}} Unsupervised learning = \textcolor{myblue}{finding patterns without labels}!
	\end{frame}
	
	\begin{frame}{T33: PCA - Village Example}
		\fontsize{6.5}{7.5}\selectfont
		\textbf{TOPIC: PCA - Village Layout Example}
		
		\vspace{0.1cm}
		\textbf{SCENARIO:}
		\begin{itemize}
			\item Plot position of houses in village
			\item Original coordinates: latitude, longitude, height
		\end{itemize}
		
		\vspace{0.1cm}
		\textbf{PCA TRANSFORMATION:}
		\begin{itemize}
			\item New framework aligns with best discrimination
			\item Village along straight road on flat land
			\item First vector: parallel to road (most spacing)
			\item Second vector: orthogonal to first
			\item Second vector: horizontal, distance from road
			\item Third vector: vertical
		\end{itemize}
		
		\vspace{0.1cm}
		\textbf{RESULT:}
		\begin{itemize}
			\item Coordinates in new frame calculated from original
			\item View from above: most information
			\item Maps standardly do this
			\item Most informative presentation of layout
		\end{itemize}
		
		\vspace{0.1cm}
		\textbf{KEY INSIGHT:}
		\begin{itemize}
			\item PCA yields what we would wish
			\item Dimensionality reduction
			\item Preserves most significant information
		\end{itemize}
		
		\textcolor{myred}{\textbf{EXAM TRAP:}} PCA finds \textcolor{myblue}{most informative dimensions} for discrimination!
	\end{frame}
	
	\begin{frame}{T34: PCA - Stylometry Example}
		\fontsize{6.5}{7.5}\selectfont
		\textbf{TOPIC: PCA - Stylometry Example (Burrows, 1982)}
		
		\vspace{0.1cm}
		\textbf{QUESTION:}
		\begin{itemize}
			\item Was "The Memoirs of a Lady of Quality" written by Smollett?
			\item Published 1751 in Smollett's Adventures of Peregrine Pickle
		\end{itemize}
		
		\vspace{0.1cm}
		\textbf{METHOD:}
		\begin{enumerate}
			\item Count occurrences of all words in corpus
			\item Order by overall frequency
			\item Select most common 50 words
			\item Count occurrences in each text
			\item Divide by total word count
			\item 50 fractional values for each text
		\end{enumerate}
		
		\vspace{0.1cm}
		\textbf{PCA APPLICATION:}
		\begin{itemize}
			\item 16 texts in 50-dimensional space
			\item Coordinates proportional to relative frequencies
			\item Texts close together: similar patterns
			\item PCA finds best 2D projection
		\end{itemize}
		
		\vspace{0.1cm}
		\textbf{RESULT:}
		\begin{itemize}
			\item Female authors clustered together
			\item "Memoirs" close to Smollett
			\item Fairly decisive positive verdict
			\item Hard to explain what similarities mean
		\end{itemize}
		
		\textcolor{myred}{\textbf{EXAM TRAP:}} PCA clusters similar texts \textcolor{myblue}{without explicit labels}!
	\end{frame}
	
	\begin{frame}{T35: Risks of Inscrutability}
		\fontsize{6.5}{7.5}\selectfont
		\textbf{TOPIC: Risks of Inscrutability}
		
		\vspace{0.1cm}
		\textbf{1. HIDDEN BIAS:}
		\begin{itemize}
			\item Bias in training data labels learned and perpetuated
			\item Presence hidden and hard to eradicate
			\item Example: Hiring algorithm favoring men
			\item Proxy variables (subjects, hobbies, careers) correlate with gender
			\item Removing explicit gender insufficient
		\end{itemize}
		
		\vspace{0.1cm}
		\textbf{2. PRIVACY RISKS:}
		\begin{itemize}
			\item Complex information entangled in network
			\item Holds clues to personal characteristics
			\item Example: Facebook Likes predict sexual orientation, ethnicity, etc.
			\item Risks of manipulation
			\item Targeted advertising
		\end{itemize}
		
		\vspace{0.1cm}
		\textbf{3. DIFFICULTY OF ERADICATION:}
		\begin{itemize}
			\item Hard to identify source of hidden biases
			\item Deep network provides no explanation
			\item True explanation hidden in vast web of weights
			\item Billions of weights
			\item Cannot be rendered humanly comprehensible
		\end{itemize}
		
		\textcolor{myred}{\textbf{EXAM TRAP:}} Inscrutability = \textcolor{myblue}{hidden bias, privacy risks, hard to eradicate}!
	\end{frame}
	
	\begin{frame}{T36: Risks of Over-Reliance}
		\fontsize{6.5}{7.5}\selectfont
		\textbf{TOPIC: Risks of Over-Reliance}
		
		\vspace{0.1cm}
		\textbf{DEFINITION:}
		\begin{itemize}
			\item Depending on system whose mode of operation seems mysterious
			\item Appears impressively expert
			\item Far faster and more confident than we could be
		\end{itemize}
		
		\vspace{0.1cm}
		\textbf{CONSEQUENCES:}
		\begin{itemize}
			\item Undermines autonomy as responsible individuals
			\item Neglect to develop or apply own judgment
			\item Fail to realize when system is wrong
		\end{itemize}
		
		\vspace{0.1cm}
		\textbf{ADVERSARIAL EXAMPLES:}
		\begin{itemize}
			\item Deep learning vulnerable to deliberately designed examples
			\item Two images differ imperceptibly to humans
			\item Categorized quite differently by system
			\item Very confidently
			\item Sticker on road sign causes misidentification
		\end{itemize}
		
		\vspace{0.1cm}
		\textbf{EXACERBATED BY HYPE:}
		\begin{itemize}
			\item Excitement about ChatGPT
			\item Supposed general intelligence
			\item Software development example
			\item Plausible but incorrect code
		\end{itemize}
		
		\textcolor{myred}{\textbf{EXAM TRAP:}} Over-reliance = \textcolor{myblue}{automation bias - complacency and reduced vigilance}!
	\end{frame}
	
	\begin{frame}{T37: Automation Bias}
		\fontsize{6.5}{7.5}\selectfont
		\textbf{TOPIC: Automation Bias}
		
		\vspace{0.1cm}
		\textbf{DEFINITION:}
		\begin{itemize}
			\item Tendency of humans to over-rely on automated systems
			\item Leads to complacency and reduced vigilance
		\end{itemize}
		
		\vspace{0.1cm}
		\textbf{CONSEQUENCES:}
		\begin{itemize}
			\item Errors of oversight
			\item Especially problematic in critical situations
			\item Skill atrophy
			\item Reduced decision-making autonomy
			\item Risk to safety and well-being
		\end{itemize}
		
		\vspace{0.1cm}
		\textbf{SKILL ATROPHY:}
		\begin{itemize}
			\item Prolonged reliance on automation
			\item Decay in essential skills
			\item Operators become passive observers
			\item From active decision makers to passive
		\end{itemize}
		
		\vspace{0.1cm}
		\textbf{MITIGATION:}
		\begin{itemize}
			\item Inform users of risks
			\item Ensure understanding of limitations
			\item Design mechanisms requiring periodic human input
			\item Keep operators engaged and alert
		\end{itemize}
		
		\textcolor{myred}{\textbf{EXAM TRAP:}} Automation bias = \textcolor{myblue}{over-reliance leading to complacency}!
	\end{frame}
	
	\begin{frame}{T38: Risks to Individuals}
		\fontsize{6.5}{7.5}\selectfont
		\textbf{TOPIC: Risks to Individuals}
		
		\vspace{0.1cm}
		\textbf{SOURCES OF RISK:}
		\begin{itemize}
			\item Organizations/institutions with whom they interact
			\item Their own behavior
			\item AI systems contributing to decision making
		\end{itemize}
		
		\vspace{0.1cm}
		\textbf{TYPES OF RISK:}
		\begin{enumerate}
			\item \textcolor{myblue}{Faulty Decisions:}
			\begin{itemize}
				\item Poorly designed AI systems
				\item Trained using inappropriate data
				\item Biased or unfair decision making
			\end{itemize}
			\item \textcolor{myblue}{Automation Bias:}
			\begin{itemize}
				\item Over-reliance on automated systems
				\item Complacency and reduced vigilance
				\item Reduced decision-making autonomy
				\item Risk to safety and well-being
			\end{itemize}
			\item \textcolor{myblue}{Manipulation:}
			\begin{itemize}
				\item Algorithms drawing on large amounts of information
				\item Predictions about future actions
				\item Influence behavior
				\item Advertisers or interest groups
			\end{itemize}
			\item \textcolor{myblue}{Privacy Threats:}
			\begin{itemize}
				\item Inferences made from information
				\item Personal characteristics exposed
			\end{itemize}
		\end{enumerate}
		
		\textcolor{myred}{\textbf{EXAM TRAP:}} Individual risks = \textcolor{myblue}{faulty decisions, automation bias, manipulation, privacy}!
	\end{frame}
	
	\begin{frame}{T39: Risks to Organizations}
		\fontsize{6.5}{7.5}\selectfont
		\textbf{TOPIC: Risks to Organizations}
		
		\vspace{0.1cm}
		\textbf{TYPES OF RISK:}
		\begin{enumerate}
			\item \textcolor{myblue}{Operational Risk:}
			\begin{itemize}
				\item Over-reliance on AI
				\item False confidence
				\item Bad decision making
				\item Customer dissatisfaction
				\item Commercial loss
			\end{itemize}
			\item \textcolor{myblue}{Reputational Risk:}
			\begin{itemize}
				\item AI practices hurting individuals or groups
				\item Poorly designed AI algorithm
				\item Lack of explainability
				\item Privacy breaches
				\item Biased results
			\end{itemize}
			\item \textcolor{myblue}{Regulatory Risk:}
			\begin{itemize}
				\item AI-specific regulations still emerging
				\item Existing privacy laws apply
				\item Model governance regulations apply
				\item Potential fines and sanctions
			\end{itemize}
			\item \textcolor{myblue}{Legal Risk:}
			\begin{itemize}
				\item Lawsuits
				\item Liability
				\item Damages
			\end{itemize}
		\end{enumerate}
		
		\textcolor{myred}{\textbf{EXAM TRAP:}} Organizational risks = \textcolor{myblue}{operational, reputational, regulatory, legal}!
	\end{frame}
	
	\begin{frame}{T40: Risks to Society}
		\fontsize{6.5}{7.5}\selectfont
		\textbf{TOPIC: Risks to Society}
		
		\vspace{0.1cm}
		\textbf{TYPES OF RISK:}
		\begin{enumerate}
			\item \textcolor{myblue}{Economic Displacement:}
			\begin{itemize}
				\item Job losses
				\item AI-driven applications taking place of human workers
				\item Widening wealth gap
				\item Skills gap
			\end{itemize}
			\item \textcolor{myblue}{Global Inequality:}
			\begin{itemize}
				\item Exacerbation between countries
				\item Wealthier, more industrialized countries advantaged
				\item Less well-positioned countries disadvantaged
			\end{itemize}
			\item \textcolor{myblue}{Misinformation:}
			\begin{itemize}
				\item Creation and spread of misinformation
				\item Deep fakes
				\item Mislead or influence public opinion
			\end{itemize}
			\item \textcolor{myblue}{Existential Risk:}
			\begin{itemize}
				\item Fear of superintelligence
				\item Autonomous weapons systems
				\item Decisions against interests of creators
				\item No consensus on nature or extent
			\end{itemize}
		\end{enumerate}
		
		\textcolor{myred}{\textbf{EXAM TRAP:}} Societal risks = \textcolor{myblue}{displacement, inequality, misinformation, existential}!
	\end{frame}
	
	\begin{frame}{T41: Over-Reliance and False Confidence}
		\fontsize{6.5}{7.5}\selectfont
		\textbf{TOPIC: Over-Reliance and False Confidence}
		
		\vspace{0.1cm}
		\textbf{THE PROBLEM:}
		\begin{itemize}
			\item Apparent expertise of AI systems
			\item Far faster and more confident than humans
			\item More sophisticated processing
			\item More comprehensive data
		\end{itemize}
		
		\vspace{0.1cm}
		\textbf{CONSEQUENCES:}
		\begin{itemize}
			\item Undermines autonomy
			\item Neglect to develop own judgment
			\item Fail to realize when system is wrong
			\item False confidence in outputs
		\end{itemize}
		
		\vspace{0.1cm}
		\textbf{SOFTWARE DEVELOPMENT EXAMPLE:}
		\begin{itemize}
			\item ChatGPT impressive in programming exercises
			\item Temptation to rely on it for code production
			\item Economize on software engineering teams
			\item But: mimicking textual patterns
			\item Not rigorously analyzing internal model
		\end{itemize}
		
		\vspace{0.1cm}
		\textbf{RISKS:}
		\begin{itemize}
			\item Non-standard tasks: incorrect code
			\item Plausible-looking code
			\item Works well in most cases
			\item Errors especially hard to identify
			\item False confidence
		\end{itemize}
		
		\textcolor{myred}{\textbf{EXAM TRAP:}} False confidence = \textcolor{myblue}{plausible but incorrect outputs}!
	\end{frame}
	
	\begin{frame}{T42: Hallucination in Large Language Models}
		\fontsize{6.5}{7.5}\selectfont
		\textbf{TOPIC: Hallucination in Large Language Models}
		
		\vspace{0.1cm}
		\textbf{DEFINITION:}
		\begin{itemize}
			\item Generate textual information completely false
			\item Yet often plausible
			\item Inevitable limitation
		\end{itemize}
		
		\vspace{0.1cm}
		\textbf{WHY IT HAPPENS:}
		\begin{itemize}
			\item LLMs based on statistical patterns
			\item Not on understanding of truth
			\item No knowledge of the world
			\item Probabilistic "guesses"
			\item Designed for plausible responses
		\end{itemize}
		
		\vspace{0.1cm}
		\textbf{EXAMPLE:}
		\begin{itemize}
			\item Plausible but false information
			\item Hard to detect
			\item Spread at scale
			\item Physical safety risks
			\item Legal risks
		\end{itemize}
		
		\vspace{0.1cm}
		\textbf{INNATE LIMITATION:}
		\begin{itemize}
			\item Hallucination is inevitable
			\item Inherent to LLMs
			\item May avoid by combining with other systems
			\item Cannot eliminate entirely
		\end{itemize}
		
		\textcolor{myred}{\textbf{EXAM TRAP:}} Hallucination = \textcolor{myblue}{false but plausible information}!
	\end{frame}
	
	\begin{frame}{T43: Adversarial Examples}
		\fontsize{6.5}{7.5}\selectfont
		\textbf{TOPIC: Adversarial Examples}
		
		\vspace{0.1cm}
		\textbf{DEFINITION:}
		\begin{itemize}
			\item Deliberately designed inputs
			\item Cause deep learning models to make errors
			\item Often imperceptible to humans
		\end{itemize}
		
		\vspace{0.1cm}
		\textbf{EXAMPLE:}
		\begin{itemize}
			\item Two images differ imperceptibly
			\item Categorized quite differently by system
			\item System very confident
			\item Sticker on road sign
			\item Stop sign misidentified as speed limit sign
		\end{itemize}
		
		\vspace{0.1cm}
		\textbf{IMPLICATIONS:}
		\begin{itemize}
			\item Security vulnerability
			\item Safety risk
			\item Potential for malicious exploitation
			\item Challenges robustness assumptions
		\end{itemize}
		
		\vspace{0.1cm}
		\textbf{WHY IT MATTERS:}
		\begin{itemize}
			\item Neural networks thought robust to noise
			\item But vulnerable to deliberate attacks
			\item Over-reliance especially dangerous
			\item Need for awareness and testing
		\end{itemize}
		
		\textcolor{myred}{\textbf{EXAM TRAP:}} Adversarial examples = \textcolor{myblue}{imperceptible changes cause major errors}!
	\end{frame}
	
	\begin{frame}{T44: Existential Risk}
		\fontsize{6.5}{7.5}\selectfont
		\textbf{TOPIC: Existential Risk from AI}
		
		\vspace{0.1cm}
		\textbf{DEFINITION:}
		\begin{itemize}
			\item Fear of "existential" risk to humanity
			\item Especially superintelligence
			\item AI system making decisions autonomously
			\item Against interests, values, priorities of creators or humanity
		\end{itemize}
		
		\vspace{0.1cm}
		\textbf{KEY CONCERNS:}
		\begin{itemize}
			\item \textcolor{myblue}{Superintelligence:} Beyond human-level intelligence
			\item \textcolor{myblue}{AI Alignment:} Ensuring means and goals aligned with human values
			\item \textcolor{myblue}{AI Arms Race:} Competition without adequate safety measures
			\item \textcolor{myblue}{Autonomous Weapons:} Decisions without human control
		\end{itemize}
		
		\vspace{0.1cm}
		\textbf{NO CONSENSUS:}
		\begin{itemize}
			\item Nature or extent of existential risk
			\item Not universally accepted concern
			\item Some argue current AI is "narrow"
			\item Far from superintelligent
			\item Regulatory mechanisms will develop
		\end{itemize}
		
		\textcolor{myred}{\textbf{EXAM TRAP:}} Existential risk = \textcolor{myblue}{no consensus on nature or extent}!
	\end{frame}
	
	\begin{frame}{T45: Narrow versus General Intelligence}
		\fontsize{6.5}{7.5}\selectfont
		\textbf{TOPIC: Narrow versus General Intelligence}
		
		\vspace{0.1cm}
		\textbf{NARROW AI (WEAK AI):}
		\begin{itemize}
			\item Designed for specific tasks
			\item Outperforms humans in narrow domains
			\item Examples: chess, image recognition, language translation
			\item All current practical AI
		\end{itemize}
		
		\vspace{0.1cm}
		\textbf{GENERAL AI (STRONG AI / AGI):}
		\begin{itemize}
			\item Human-level general intelligence
			\item Can handle any intellectual task
			\item Not yet achieved
			\item Source of existential risk concerns
		\end{itemize}
		
		\vspace{0.1cm}
		\textbf{KEY DISTINCTION:}
		\begin{itemize}
			\item Specific intelligence: good at one task
			\item General intelligence: good at many tasks
			\item Current AI: narrow, not general
			\item Turing Test: conceptual existence proof
		\end{itemize}
		
		\vspace{0.1cm}
		\textbf{IMPLICATIONS:}
		\begin{itemize}
			\item Machines can be intelligent at specific problems
			\item Without general competence
			\item Less threatening than AGI
			\item Focus of practical AI applications
		\end{itemize}
		
		\textcolor{myred}{\textbf{EXAM TRAP:}} Most AI is \textcolor{myblue}{narrow}, not general!
	\end{frame}
	
	\begin{frame}{T46: Statistical Pattern Recognition}
		\fontsize{6.5}{7.5}\selectfont
		\textbf{TOPIC: Statistical Pattern Recognition}
		
		\vspace{0.1cm}
		\textbf{DEFINITION:}
		\begin{itemize}
			\item How neural networks work
			\item Based on statistical patterns in data
			\item Not on understanding or explicit reasoning
		\end{itemize}
		
		\vspace{0.1cm}
		\textbf{KEY CHARACTERISTICS:}
		\begin{itemize}
			\item Learn patterns from training data
			\item Represent patterns implicitly in weights
			\item Make predictions based on patterns
			\item No explicit model of domain
		\end{itemize}
		
		\vspace{0.1cm}
		\textbf{ANALOGY TO HUMAN COGNITION:}
		\begin{itemize}
			\item Similar to unconscious pattern-recognition
			\item Unlike explicit calculation
			\item We recognize things without knowing how
			\item Hard to articulate our own processes
		\end{itemize}
		
		\vspace{0.1cm}
		\textbf{IMPLICATIONS:}
		\begin{itemize}
			\item Inscrutability
			\item Hidden biases
			\item Limited understanding
			\item "Stochastic parrot" nature
		\end{itemize}
		
		\textcolor{myred}{\textbf{EXAM TRAP:}} Neural networks = \textcolor{myblue}{statistical pattern recognition}, not understanding!
	\end{frame}
	
	\begin{frame}{T47: Training Data and Bias}
		\fontsize{6.5}{7.5}\selectfont
		\textbf{TOPIC: Training Data and Bias}
		
		\vspace{0.1cm}
		\textbf{HOW BIAS ENTERS:}
		\begin{itemize}
			\item Bias in labelling of training data
			\item Historical prejudice
			\item Learned by model
			\item Perpetuated in outputs
		\end{itemize}
		
		\vspace{0.1cm}
		\textbf{HIRING EXAMPLE:}
		\begin{itemize}
			\item Company historically favored men
			\item Data about past recruitment
			\item CVs and acceptance/rejection
			\item Model exhibits bias favoring men
			\item Bias assists in matching training labels
		\end{itemize}
		
		\vspace{0.1cm}
		\textbf{PROXY VARIABLES:}
		\begin{itemize}
			\item Removing gender insufficient
			\item Other information acts as proxy
			\item Subjects at school
			\item Sports and hobbies
			\item Careers
			\item Strongly correlated with gender
		\end{itemize}
		
		\vspace{0.1cm}
		\textbf{LOAN EXAMPLE:}
		\begin{itemize}
			\item Historically racially biased lending
			\item Segregated city areas
			\item Removing race from data
			\item Postal code correlates with race
			\item Perpetuates prejudice
		\end{itemize}
		
		\textcolor{myred}{\textbf{EXAM TRAP:}} Removing explicit bias \textcolor{myblue}{insufficient} due to proxies!
	\end{frame}
	
	\begin{frame}{T48: Privacy Risks from AI}
		\fontsize{6.5}{7.5}\selectfont
		\textbf{TOPIC: Privacy Risks from AI}
		
		\vspace{0.1cm}
		\textbf{HOW PRIVACY IS THREATENED:}
		\begin{itemize}
			\item Deep implicit entanglement of information
			\item Network holds clues to personal characteristics
			\item Characteristics we prefer to keep secret
		\end{itemize}
		
		\vspace{0.1cm}
		\textbf{FACEBOOK LIKES EXAMPLE:}
		\begin{itemize}
			\item Used to predict sensitive attributes
			\item Sexual orientation
			\item Ethnicity
			\item Religious and political views
			\item Personality traits
			\item Intelligence
			\item Happiness
			\item Use of addictive substances
			\item Parental separation
			\item Age and gender
		\end{itemize}
		
		\vspace{0.1cm}
		\textbf{RISKS OF MANIPULATION:}
		\begin{itemize}
			\item Commercial sphere
			\item Political sphere
			\item Identify personal characteristics
			\item Target advertising
			\item Influence behavior
		\end{itemize}
		
		\textcolor{myred}{\textbf{EXAM TRAP:}} AI can \textcolor{myblue}{infer sensitive personal characteristics} from data!
	\end{frame}
	
	\begin{frame}{T49: Explainable AI Challenge}
		\fontsize{6.5}{7.5}\selectfont
		\textbf{TOPIC: The Explainable AI Challenge}
		
		\vspace{0.1cm}
		\textbf{THE PROBLEM:}
		\begin{itemize}
			\item Deep network models cannot explain decisions
			\item True explanation hidden in vast web of weights
			\item Billions of weights
			\item Cannot be rendered humanly comprehensible
		\end{itemize}
		
		\vspace{0.1cm}
		\textbf{CONTRAST WITH EXPERT SYSTEMS:}
		\begin{itemize}
			\item Expert systems: explicit rules
			\item Can explain reasoning
			\item Deep networks: no explanation
			\item No way to identify source of bias
		\end{itemize}
		
		\vspace{0.1cm}
		\textbf{CHALLENGE:}
		\begin{itemize}
			\item Achieving explainable AI
			\item While retaining power of machine learning
			\item "Black box" to "glass box"
			\item Major research area
		\end{itemize}
		
		\vspace{0.1cm}
		\textbf{IMPLICATIONS:}
		\begin{itemize}
			\item Huge reputational risk
			\item Regulatory challenges
			\item Difficulty eradicating bias
			\item Trust issues
		\end{itemize}
		
		\textcolor{myred}{\textbf{EXAM TRAP:}} Explainable AI = \textcolor{myblue}{making black boxes transparent}!
	\end{frame}
	
	\begin{frame}{T50: Course Navigation - Tools and Techniques}
		\fontsize{6.5}{7.5}\selectfont
		\textbf{TOPIC: Course Navigation - Tools and Techniques}
		
		\vspace{0.1cm}
		\textbf{FOCUS:}
		\begin{itemize}
			\item Mastering the AI toolbox itself
			\item Core building blocks of AI
			\item Categorization of ML algorithms
		\end{itemize}
		
		\vspace{0.1cm}
		\textbf{TOPICS COVERED:}
		\begin{enumerate}
			\item \textcolor{myblue}{Machine Learning Algorithms:}
			\begin{itemize}
				\item Supervised learning
				\item Unsupervised learning
				\item Reinforcement learning
			\end{itemize}
			\item \textcolor{myblue}{Data Preparation:}
			\begin{itemize}
				\item Transforming raw information
				\item Insights fit for modeling
			\end{itemize}
			\item \textcolor{myblue}{Feature Selection:}
			\begin{itemize}
				\item Choosing relevant features
				\item Improving model performance
			\end{itemize}
			\item \textcolor{myblue}{Visualization:}
			\begin{itemize}
				\item Understanding data
				\item Communicating insights
			\end{itemize}
			\item \textcolor{myblue}{Explainability:}
			\begin{itemize}
				\item Balancing accuracy and explainability
				\item Vital for business decision-making
			\end{itemize}
		\end{enumerate}
		
		\textcolor{myred}{\textbf{KEY POINT:}} Mastering tools is \textcolor{myblue}{essential foundation}!
	\end{frame}
	
	\begin{frame}{T51: Course Navigation - Risks and Risk Factors}
		\fontsize{6.5}{7.5}\selectfont
		\textbf{TOPIC: Course Navigation - Risks and Risk Factors}
		
		\vspace{0.1cm}
		\textbf{FOCUS:}
		\begin{itemize}
			\item Potential unintended consequences of AI
			\item Broader perspective
		\end{itemize}
		
		\vspace{0.1cm}
		\textbf{TOPICS COVERED:}
		\begin{enumerate}
			\item \textcolor{myblue}{Bias Amplification:}
			\begin{itemize}
				\item Existing biases amplified
				\item New threats to autonomy
				\item Safety concerns
			\end{itemize}
			\item \textcolor{myblue}{Mitigation Strategies:}
			\begin{itemize}
				\item Fairness
				\item Transparency
				\item Complex error translation
			\end{itemize}
			\item \textcolor{myblue}{Real-World Harm:}
			\begin{itemize}
				\item Errors translating into harm
				\item Reputational risks
				\item Public mistrust
			\end{itemize}
			\item \textcolor{myblue}{AI Governance:}
			\begin{itemize}
				\item Proactive governance
				\item Building trustworthy organization
			\end{itemize}
		\end{enumerate}
		
		\textcolor{myred}{\textbf{KEY POINT:}} Understanding risks is \textcolor{myblue}{essential for responsible AI}!
	\end{frame}
	
	\begin{frame}{T52: Course Navigation - Ethical and Responsible AI}
		\fontsize{6.5}{7.5}\selectfont
		\textbf{TOPIC: Course Navigation - Ethical and Responsible AI}
		
		\vspace{0.1cm}
		\textbf{FOCUS:}
		\begin{itemize}
			\item Moving beyond risk mitigation
			\item Principled implementation
		\end{itemize}
		
		\vspace{0.1cm}
		\textbf{TOPICS COVERED:}
		\begin{enumerate}
			\item \textcolor{myblue}{Ethical Frameworks:}
			\begin{itemize}
				\item Consequentialism
				\item Deontology
				\item Virtue ethics
				\item Applied to real-world scenarios
			\end{itemize}
			\item \textcolor{myblue}{Legal Landscape:}
			\begin{itemize}
				\item EU AI Act
				\item US regulatory initiatives
				\item Evolving regulations
			\end{itemize}
			\item \textcolor{myblue}{Proactive Approach:}
			\begin{itemize}
				\item Not merely reactive
				\item Key leadership quality
				\item Rules still being shaped
			\end{itemize}
		\end{enumerate}
		
		\textcolor{myred}{\textbf{KEY POINT:}} Ethics and regulation are \textcolor{myblue}{critical for responsible AI}!
	\end{frame}
	
	\begin{frame}{T53: Course Navigation - Governance}
		\fontsize{6.5}{7.5}\selectfont
		\textbf{TOPIC: Course Navigation - Governance}
		
		\vspace{0.1cm}
		\textbf{FOCUS:}
		\begin{itemize}
			\item Critical lever for successful implementation
			\item Not bureaucratic burden
		\end{itemize}
		
		\vspace{0.1cm}
		\textbf{TOPICS COVERED:}
		\begin{enumerate}
			\item \textcolor{myblue}{Data Governance:}
			\begin{itemize}
				\item Robust data strategy
				\item Ensuring data quality
				\item Data provenance
			\end{itemize}
			\item \textcolor{myblue}{Model Governance:}
			\begin{itemize}
				\item Clear processes for development
				\item Testing
				\item Ongoing monitoring
			\end{itemize}
			\item \textcolor{myblue}{Outcomes:}
			\begin{itemize}
				\item Reliable insights
				\item Avoiding unexpected liabilities
				\item Avoiding operational weaknesses
			\end{itemize}
		\end{enumerate}
		
		\textcolor{myred}{\textbf{KEY POINT:}} Governance ensures \textcolor{myblue}{reliable and responsible AI}!
	\end{frame}
	
	\begin{frame}{T54: Course Navigation - Case Studies}
		\fontsize{6.5}{7.5}\selectfont
		\textbf{TOPIC: Course Navigation - Case Studies and Practitioner Perspectives}
		
		\vspace{0.1cm}
		\textbf{FOCUS:}
		\begin{itemize}
			\item Bringing concepts to life
			\item Practical applications
		\end{itemize}
		
		\vspace{0.1cm}
		\textbf{TOPICS COVERED:}
		\begin{enumerate}
			\item \textcolor{myblue}{Case Studies:}
			\begin{itemize}
				\item Specific applications in finance
				\item Risk management examples
				\item Real-world challenges
			\end{itemize}
			\item \textcolor{myblue}{Practitioner Perspectives:}
			\begin{itemize}
				\item How ML transforms risk modeling
				\item Practical strategies
				\item Ensuring explainability
				\item Robustness in high-stakes environments
			\end{itemize}
		\end{enumerate}
		
		\textcolor{myred}{\textbf{KEY POINT:}} Case studies provide \textcolor{myblue}{practical context}!
	\end{frame}
	
	\begin{frame}{T55: Course Goal - Business Leadership}
		\fontsize{6.5}{7.5}\selectfont
		\textbf{TOPIC: Course Goal - Business Leadership}
		
		\vspace{0.1cm}
		\textbf{OBJECTIVE:}
		\begin{itemize}
			\item Not to become a data scientist
			\item Rather a business leader
			\item Speak the language of data scientists
			\item Bridge gap between technical possibilities and business needs
		\end{itemize}
		
		\vspace{0.1cm}
		\textbf{KEY SKILLS:}
		\begin{enumerate}
			\item \textcolor{myblue}{Understanding AI:}
			\begin{itemize}
				\item Core concepts
				\item Capabilities and limitations
			\end{itemize}
			\item \textcolor{myblue}{Risk Management:}
			\begin{itemize}
				\item Identifying risks
				\item Mitigating risks
				\item Governing AI
			\end{itemize}
			\item \textcolor{myblue}{Ethics:}
			\begin{itemize}
				\item Ethical frameworks
				\item Responsible implementation
			\end{itemize}
			\item \textcolor{myblue}{Communication:}
			\begin{itemize}
				\item Explaining AI to stakeholders
				\item Translating technical to business
			\end{itemize}
		\end{enumerate}
		
		\vspace{0.1cm}
		\textbf{ULTIMATE GOAL:}
		\begin{itemize}
			\item Unlock transformative potential of AI
			\item Sustainable
			\item Responsible
			\item Profitable
		\end{itemize}
		
		\textcolor{myred}{\textbf{KEY POINT:}} Business leaders must \textcolor{myblue}{bridge technical and business worlds}!
	\end{frame}
	
	\begin{frame}{T56: Historical Perspective - Babbage}
		\fontsize{6.5}{7.5}\selectfont
		\textbf{TOPIC: Charles Babbage and Early Computing}
		
		\vspace{0.1cm}
		\textbf{KEY CONTRIBUTIONS:}
		\begin{itemize}
			\item Early 19th century
			\item Difference engine: mechanical, gear-driven
			\item Could replace error-prone human calculation
			\item Efficient production of mathematical tables
		\end{itemize}
		
		\vspace{0.1cm}
		\textbf{ANALYTICAL ENGINE:}
		\begin{itemize}
			\item More ambitious design
			\item Widely seen as anticipating digital computer
			\item Never completed
			\item Huge feat of engineering
		\end{itemize}
		
		\vspace{0.1cm}
		\textbf{SIGNIFICANCE:}
		\begin{itemize}
			\item First practical vision of computing
			\item Inspired future developments
			\item Showed possibility of automated calculation
		\end{itemize}
		
		\textcolor{myred}{\textbf{EXAM TRAP:}} Babbage = \textcolor{myblue}{early vision of computing}!
	\end{frame}
	
	\begin{frame}{T57: Alan Turing and Computability}
		\fontsize{6.5}{7.5}\selectfont
		\textbf{TOPIC: Alan Turing and Computability}
		
		\vspace{0.1cm}
		\textbf{KEY FACTS:}
		\begin{itemize}
			\item 1936: Turing machine concept
			\item Theoretical invention
			\item Proved fundamental results about mathematics
			\item Theory of computation
		\end{itemize}
		
		\vspace{0.1cm}
		\textbf{TURING MACHINE:}
		\begin{itemize}
			\item Universal programmable computer
			\item Can follow any given recipe for calculation
			\item Not a practical device
			\item Theoretical foundation
		\end{itemize}
		
		\vspace{0.1cm}
		\textbf{BLETCHLEY PARK:}
		\begin{itemize}
			\item WWII code-breaking
			\item Mechanical methods of decrypting Enigma codes
			\item Massive impetus to electronic computer development
		\end{itemize}
		
		\vspace{0.1cm}
		\textbf{LATER WORK:}
		\begin{itemize}
			\item 1950 paper: "Computing Machinery and Intelligence"
			\item Imitation game (Turing Test)
			\item Visionary thoughts on intelligent machines
		\end{itemize}
		
		\textcolor{myred}{\textbf{EXAM TRAP:}} Turing = \textcolor{myblue}{founder of computability theory and AI}!
	\end{frame}
	
	\begin{frame}{T58: Deep Blue and Chess}
		\fontsize{6.5}{7.5}\selectfont
		\textbf{TOPIC: Deep Blue and Chess}
		
		\vspace{0.1cm}
		\textbf{KEY FACTS:}
		\begin{itemize}
			\item IBM's chess-playing computer
			\item Defeated Garry Kasparov in 1997
			\item Kasparov considered strongest human player
			\item Conspicuous success of classical AI
		\end{itemize}
		
		\vspace{0.1cm}
		\textbf{CHARACTERISTICS:}
		\begin{itemize}
			\item Specifically programmed to play chess
			\item Input from expert players
			\item Algorithms finely tuned to established theory
			\item Exhaustive search with heuristics
		\end{itemize}
		
		\vspace{0.1cm}
		\textbf{CONTRAST WITH ALPHAZERO:}
		\begin{itemize}
			\item Deep Blue: human knowledge input
			\item AlphaZero: self-taught
			\item Deep Blue: chess-specific
			\item AlphaZero: generalizable
		\end{itemize}
		
		\textcolor{myred}{\textbf{EXAM TRAP:}} Deep Blue = \textcolor{myblue}{classical AI triumph}!
	\end{frame}
	
	\begin{frame}{T59: Samuel's Checkers Program}
		\fontsize{6.5}{7.5}\selectfont
		\textbf{TOPIC: Samuel's Checkers Program}
		
		\vspace{0.1cm}
		\textbf{KEY FACTS:}
		\begin{itemize}
			\item Developed by Arthur Samuel (late 1950s)
			\item IBM researcher
			\item First substantial example of machine learning
			\item Samuel coined the term "machine learning"
		\end{itemize}
		
		\vspace{0.1cm}
		\textbf{APPROACH:}
		\begin{itemize}
			\item Limited exhaustive search to n moves
			\item Used heuristics to assess positions
			\item Material balance
			\item Mobility
			\item Advancement
			\item Kings
		\end{itemize}
		
		\vspace{0.1cm}
		\textbf{REINFORCEMENT LEARNING:}
		\begin{itemize}
			\item Weights adjusted based on game outcomes
			\item Simulated evolution
			\item Program learned to play better than designer
			\item 38 criteria devised
		\end{itemize}
		
		\textcolor{myred}{\textbf{EXAM TRAP:}} Samuel's checkers = \textcolor{myblue}{first substantial machine learning}!
	\end{frame}
	
	\begin{frame}{T60: Heuristic Search - A* Algorithm}
		\fontsize{6.5}{7.5}\selectfont
		\textbf{TOPIC: A* Algorithm}
		
		\vspace{0.1cm}
		\textbf{DEFINITION:}
		\begin{itemize}
			\item Best-known heuristic search technique
			\item Used for route-finding
			\item Finds optimal path
		\end{itemize}
		
		\vspace{0.1cm}
		\textbf{HOW IT WORKS:}
		\begin{itemize}
			\item For each location, track:
			\begin{itemize}
				\item Distance traveled so far
				\item Optimistic estimate of distance remaining
			\end{itemize}
			\item Order locations by sum of these values
			\item Always explore from location with shortest sum
			\item Continue until destination reached
		\end{itemize}
		
		\vspace{0.1cm}
		\textbf{KEY REQUIREMENT:}
		\begin{itemize}
			\item Estimate must be optimistic
			\item Cannot overestimate remaining distance
			\item Guarantees finding best route
			\item Euclidean distance suitable (crow-flies)
		\end{itemize}
		
		\vspace{0.1cm}
		\textbf{APPLICATIONS:}
		\begin{itemize}
			\item Route planning
			\item Game playing
			\item Various optimization problems
		\end{itemize}
		
		\textcolor{myred}{\textbf{EXAM TRAP:}} A* = \textcolor{myblue}{optimal heuristic search}!
	\end{frame}
	
	\begin{frame}{T61: Nim Game - Detailed Example}
		\fontsize{6.5}{7.5}\selectfont
		\textbf{TOPIC: Nim Game - Detailed Reinforcement Learning}
		
		\vspace{0.1cm}
		\textbf{GAME RULES:}
		\begin{itemize}
			\item 20 coins on table
			\item Players alternate
			\item Remove 1, 2, or 3 coins
			\item Winner removes last coin(s)
		\end{itemize}
		
		\vspace{0.1cm}
		\textbf{WINNING STRATEGY:}
		\begin{itemize}
			\item Second player has winning strategy
			\item Reduce to multiples of 4
			\item If opponent removes 1, remove 3
			\item If opponent removes 2, remove 2
			\item If opponent removes 3, remove 1
			\item Total of 4 removed per round
		\end{itemize}
		
		\vspace{0.1cm}
		\textbf{LEARNING PROCESS:}
		\begin{itemize}
			\item Initial weights: all 5
			\item Choose randomly based on weights
			\item After game: adjust weights
			\item Winner: chosen moves +1, others -1
			\item Loser: chosen moves -1, others +1
			\item Weights range: 0 to 9
		\end{itemize}
		
		\vspace{0.1cm}
		\textbf{RESULT:}
		\begin{itemize}
			\item After ~100 games: discovers winning strategy
			\item Weights reflect optimal moves
			\item 9-0-0 pattern for multiples of 4
		\end{itemize}
		
		\textcolor{myred}{\textbf{EXAM TRAP:}} Nim demonstrates \textcolor{myblue}{simple reinforcement learning}!
	\end{frame}
	
	\begin{frame}{T62: Tower of Hanoi - Recursive Solution}
		\fontsize{6.5}{7.5}\selectfont
		\textbf{TOPIC: Tower of Hanoi - Recursive Solution}
		
		\vspace{0.1cm}
		\textbf{PROBLEM:}
		\begin{itemize}
			\item Nine disks of different sizes
			\item Originally piled largest to smallest
			\item Move entire tower from left to right
			\item Only one disk at a time
			\item Never put larger disk on smaller
		\end{itemize}
		
		\vspace{0.1cm}
		\textbf{RECURSIVE SOLUTION:}
		\begin{itemize}
			\item To move pile of n disks from A to B:
			\begin{enumerate}
				\item Move pile of n-1 disks from A to C
				\item Move disk n from A to B
				\item Move pile of n-1 disks from C to B
			\end{enumerate}
			\item Each subtask reduces to simpler instance
			\item Eventually reduces to single disk movements
		\end{itemize}
		
		\vspace{0.1cm}
		\textbf{MATHEMATICAL RESULT:}
		\begin{itemize}
			\item Moving n disks takes $2^n - 1$ steps
			\item Moving 9 disks takes 511 steps
			\item Proved by mathematical induction
		\end{itemize}
		
		\vspace{0.1cm}
		\textbf{KEY INSIGHT:}
		\begin{itemize}
			\item Complex problem solved elegantly
			\item No need for lookahead or comparison
			\item Recursion provides complete solution
		\end{itemize}
		
		\textcolor{myred}{\textbf{EXAM TRAP:}} Tower of Hanoi = \textcolor{myblue}{elegant recursive solution}!
	\end{frame}
	
	\begin{frame}{T63: Noughts and Crosses - Complete Algorithm}
		\fontsize{6.5}{7.5}\selectfont
		\textbf{TOPIC: Noughts and Crosses - Complete Algorithm}
		
		\vspace{0.1cm}
		\textbf{POSITION EVALUATION:}
		\begin{itemize}
			\item Value +1: Win for player to move
			\item Value -1: Loss for player to move
			\item Value 0: Draw
		\end{itemize}
		
		\vspace{0.1cm}
		\textbf{ALGORITHM:}
		\begin{enumerate}
			\item If X has lost (opponent has line of three): value = -1
			\item If position is full: value = 0
			\item Otherwise: construct all positions with Y to move
			\item Evaluate each from Y's point of view
			\item Assign X-value = inverse of lowest Y-value
			\item When turn comes, choose move with lowest opponent value
		\end{enumerate}
		
		\vspace{0.1cm}
		\textbf{MINIMAX PRINCIPLE:}
		\begin{itemize}
			\item Presume opponent chooses best reply
			\item Choose move minimizing opponent's maximum benefit
			\item Applied recursively at every level
		\end{itemize}
		
		\vspace{0.1cm}
		\textbf{IMPLEMENTATION:}
		\begin{itemize}
			\item Fewer than 25 short lines of code
			\item Exhaustively analyzes all possible lines
			\item Creates implicit search tree
		\end{itemize}
		
		\textcolor{myred}{\textbf{EXAM TRAP:}} Minimax = \textcolor{myblue}{adversarial search principle}!
	\end{frame}
	
	\begin{frame}{T64: Heuristic Evaluation in Games}
		\fontsize{6.5}{7.5}\selectfont
		\textbf{TOPIC: Heuristic Evaluation in Games}
		
		\vspace{0.1cm}
		\textbf{PURPOSE:}
		\begin{itemize}
			\item Assess positions when exhaustive search impractical
			\item Guide choice of position to aim for
			\item Cannot guarantee optimal solution
		\end{itemize}
		
		\vspace{0.1cm}
		\textbf{CHECKERS CRITERIA (SAMUEL):}
		\begin{enumerate}
			\item Material balance
			\item Relative mobility
			\item Advancement
			\item Kings
			\item 38 criteria total
			\item 16 used at any time
		\end{enumerate}
		
		\vspace{0.1cm}
		\textbf{WEIGHTED EVALUATION:}
		\[
		v = w_1 c_1 + w_2 c_2 + ... + w_6 c_6
		\]
		
		\vspace{0.1cm}
		\textbf{WEIGHT REFINEMENT:}
		\begin{itemize}
			\item Initial weights equal
			\item Random variation creates competing strategy
			\item Play games between strategies
			\item Replace current if competing better
			\item Iterative improvement
		\end{itemize}
		
		\vspace{0.1cm}
		\textbf{KEY INSIGHT:}
		\begin{itemize}
			\item Human judgment: identify criteria
			\item Reinforcement learning: optimize weights
			\item Program can learn to beat designer
		\end{itemize}
		
		\textcolor{myred}{\textbf{EXAM TRAP:}} Heuristic weights \textcolor{myblue}{refined by reinforcement learning}!
	\end{frame}
	
	\begin{frame}{T65: Expert Systems - PROLOG}
		\fontsize{6.5}{7.5}\selectfont
		\textbf{TOPIC: Expert Systems and PROLOG}
		
		\vspace{0.1cm}
		\textbf{DEFINITION:}
		\begin{itemize}
			\item Systems designed to capture knowledge in specific domains
			\item Use logic programming
			\item Response to limitations of general AI
		\end{itemize}
		
		\vspace{0.1cm}
		\textbf{PROLOG:}
		\begin{itemize}
			\item Logic programming language
			\item Used as inferential medium
			\item Part of Japanese Fifth Generation project
		\end{itemize}
		
		\vspace{0.1cm}
		\textbf{ARCHITECTURE:}
		\begin{itemize}
			\item Knowledge base: domain-specific facts
			\item Inference engine: general reasoning
			\item Knowledge engineers: extend to new domains
		\end{itemize}
		
		\vspace{0.1cm}
		\textbf{LIMITATIONS:}
		\begin{itemize}
			\item Not easily generalized
			\item Same old problems beyond domains
			\item Brittle - fail unexpectedly
			\item Hard to extend
		\end{itemize}
		
		\textcolor{myred}{\textbf{EXAM TRAP:}} Expert systems = \textcolor{myblue}{domain-specific knowledge + general inference}!
	\end{frame}
	
	\begin{frame}{T66: Frame Problem}
		\fontsize{6.5}{7.5}\selectfont
		\textbf{TOPIC: The Frame Problem}
		
		\vspace{0.1cm}
		\textbf{DEFINITION:}
		\begin{itemize}
			\item Keeping track of what changes when action performed
			\item Also tracking what stays the same
			\item Adding to combinatorial explosion
		\end{itemize}
		
		\vspace{0.1cm}
		\textbf{THE CHALLENGE:}
		\begin{itemize}
			\item Many aspects of situation
			\item Action changes some aspects
			\item Action doesn't change others
			\item Must represent both
		\end{itemize}
		
		\vspace{0.1cm}
		\textbf{CONSEQUENCES:}
		\begin{itemize}
			\item Detailed codification adds complexity
			\item Fuel to combinatorial explosion
			\item Makes general search impractical
			\item Problem for classical AI
		\end{itemize}
		
		\vspace{0.1cm}
		\textbf{RELEVANCE:}
		\begin{itemize}
			\item Shows limits of classical AI
			\item Need for different approaches
			\item Motivates machine learning
		\end{itemize}
		
		\textcolor{myred}{\textbf{EXAM TRAP:}} Frame problem = \textcolor{myblue}{tracking what changes and what doesn't}!
	\end{frame}
	
	\begin{frame}{T67: Brittleness of AI Systems}
		\fontsize{6.5}{7.5}\selectfont
		\textbf{TOPIC: Brittleness of AI Systems}
		
		\vspace{0.1cm}
		\textbf{DEFINITION:}
		\begin{itemize}
			\item AI systems fail in unexpected ways
			\item Beyond familiar narrow boundaries
			\item Situations not explicitly foreseen
		\end{itemize}
		
		\vspace{0.1cm}
		\textbf{CHARACTERISTICS:}
		\begin{itemize}
			\item Notorious for failing unexpectedly
			\item Work well within boundaries
			\item Fail when boundaries exceeded
			\item Hard to predict failures
		\end{itemize}
		
		\vspace{0.1cm}
		\textbf{IMPLICATIONS:}
		\begin{itemize}
			\item Need for extensive testing
			\item Careful deployment
			\item Awareness of limitations
			\item Human oversight
		\end{itemize}
		
		\vspace{0.1cm}
		\textbf{RESPONSE:}
		\begin{itemize}
			\item Expert systems for specific domains
			\item Machine learning for flexibility
			\item Hybrid approaches
		\end{itemize}
		
		\textcolor{myred}{\textbf{EXAM TRAP:}} Brittleness = \textcolor{myblue}{failing unexpectedly beyond boundaries}!
	\end{frame}
	
	\begin{frame}{T68: Deep Learning Renaissance}
		\fontsize{6.5}{7.5}\selectfont
		\textbf{TOPIC: Deep Learning Renaissance}
		
		\vspace{0.1cm}
		\textbf{TIMELINE:}
		\begin{itemize}
			\item 1980s: Backpropagation rises
			\item 1990s: Enthusiasm declines
			\item 1998: LeNet-5 breakthrough
			\item 2012: AlexNet triumph
			\item 2013: Atari games
			\item 2014-2016: AlphaGo
			\item 2017: AlphaZero
			\item 2022: ChatGPT
		\end{itemize}
		
		\vspace{0.1cm}
		\textbf{ENABLING FACTORS:}
		\begin{itemize}
			\item Computational power (GPUs)
			\item Large datasets
			\item Algorithmic advances
			\item Investment and interest
		\end{itemize}
		
		\vspace{0.1cm}
		\textbf{KEY APPLICATIONS:}
		\begin{itemize}
			\item Image recognition
			\item Speech recognition
			\item Game playing
			\item Language models
		\end{itemize}
		
		\textcolor{myred}{\textbf{EXAM TRAP:}} Deep learning renaissance driven by \textcolor{myblue}{computation, data, and algorithms}!
	\end{frame}
	
	\begin{frame}{T69: ImageNet Competition}
		\fontsize{6.5}{7.5}\selectfont
		\textbf{TOPIC: ImageNet Competition}
		
		\vspace{0.1cm}
		\textbf{KEY FACTS:}
		\begin{itemize}
			\item High-profile international competition
			\item Identify images from ImageNet collection
			\item Built by Fei Fei Li (started 2006)
			\item 14.2 million images
			\item Nearly 22,000 categories
		\end{itemize}
		
		\vspace{0.1cm}
		\textbf{ALEXNET WIN (2012):}
		\begin{itemize}
			\item Error rate: 15.3\%
			\item Previous year: 25.8\%
			\item Nearly 100 times bigger than LeNet-5
			\item Almost 900,000 neurons
		\end{itemize}
		
		\vspace{0.1cm}
		\textbf{IMPACT:}
		\begin{itemize}
			\item Startling vindication of deep learning
			\item Demonstrated amazing power
			\item Sparked renewed interest
			\item Changed computer vision field
		\end{itemize}
		
		\textcolor{myred}{\textbf{EXAM TRAP:}} AlexNet (2012) = \textcolor{myblue}{breakthrough in image recognition}!
	\end{frame}
	
	\begin{frame}{T70: GPU Revolution}
		\fontsize{6.5}{7.5}\selectfont
		\textbf{TOPIC: GPU Revolution}
		
		\vspace{0.1cm}
		\textbf{DEFINITION:}
		\begin{itemize}
			\item Graphics Processing Units
			\item Developed for animated computer games
			\item Process image pixels in parallel
		\end{itemize}
		
		\vspace{0.1cm}
		\textbf{CHARACTERISTICS:}
		\begin{itemize}
			\item Up to 10,000 cores
			\item Massive parallel processing
			\item Perfect for matrix operations
			\item Ideal for neural networks
		\end{itemize}
		
		\vspace{0.1cm}
		\textbf{IMPACT ON AI:}
		\begin{itemize}
			\item Enabled deep learning
			\item Made AlexNet possible
			\item Accelerated training
			\item Democratized AI research
		\end{itemize}
		
		\vspace{0.1cm}
		\textbf{BROADER SIGNIFICANCE:}
		\begin{itemize}
			\item Gaming industry drove innovation
			\item Unexpected benefit for AI
			\item Example of technology transfer
		\end{itemize}
		
		\textcolor{myred}{\textbf{EXAM TRAP:}} GPUs enabled \textcolor{myblue}{deep learning revolution}!
	\end{frame}
	
	\begin{frame}{T71: AlphaGo - Significance}
		\fontsize{6.5}{7.5}\selectfont
		\textbf{TOPIC: AlphaGo - Significance}
		
		\vspace{0.1cm}
		\textbf{GO GAME:}
		\begin{itemize}
			\item Ancient Chinese board game
			\item More complex than chess
			\item Average branching factor ~250
			\item Considered too subtle for computers
		\end{itemize}
		
		\vspace{0.1cm}
		\textbf{ALPHAGO ACHIEVEMENTS:}
		\begin{itemize}
			\item Defeated Fan Hui 5-0 (2014)
			\item Defeated Lee Sedol 4-1 (2015)
			\item World champion defeated
			\item Historic milestone
		\end{itemize}
		
		\vspace{0.1cm}
		\textbf{TECHNIQUES:}
		\begin{itemize}
			\item Deep convolutional neural networks
			\item Reinforcement learning
			\item Monte Carlo tree search
			\item Combination of approaches
		\end{itemize}
		
		\vspace{0.1cm}
		\textbf{SIGNIFICANCE:}
		\begin{itemize}
			\item Showed AI can master complex games
			\item Beyond human capability
			\item Not specifically programmed
			\item Learned from experience
		\end{itemize}
		
		\textcolor{myred}{\textbf{EXAM TRAP:}} AlphaGo = \textcolor{myblue}{mastering complex game of Go}!
	\end{frame}
	
	\begin{frame}{T72: AlphaZero - Generalization}
		\fontsize{6.5}{7.5}\selectfont
		\textbf{TOPIC: AlphaZero - Generalization}
		
		\vspace{0.1cm}
		\textbf{ADVANCE OVER ALPHAGO:}
		\begin{itemize}
			\item More generalized program
			\item Capable of teaching itself new games
			\item Go and chess
			\item Not just one game
		\end{itemize}
		
		\vspace{0.1cm}
		\textbf{WHAT IT WAS GIVEN:}
		\begin{itemize}
			\item Starting position
			\item Moves of all pieces
			\item Various rules
			\item Win, draw, loss identification
			\item Efficient recursive adversarial tree search
		\end{itemize}
		
		\vspace{0.1cm}
		\textbf{HOW IT LEARNED:}
		\begin{itemize}
			\item Competing against itself
			\item No human expert input
			\item No game databases
			\item Self-training for few hours
			\item Defeated Stockfish (world champion chess program)
		\end{itemize}
		
		\vspace{0.1cm}
		\textbf{SIGNIFICANCE:}
		\begin{itemize}
			\item Massive progress toward AGI
			\item Self-taught skills
			\item Beyond human ingenuity
			\item Potentially frightening
		\end{itemize}
		
		\textcolor{myred}{\textbf{EXAM TRAP:}} AlphaZero = \textcolor{myblue}{self-teaching through self-play}!
	\end{frame}
	
	\begin{frame}{T73: Stochastic Parrot}
		\fontsize{6.5}{7.5}\selectfont
		\textbf{TOPIC: Stochastic Parrot}
		
		\vspace{0.1cm}
		\textbf{DEFINITION:}
		\begin{itemize}
			\item Term coined by Emily M. Bender (2021)
			\item Describes LLMs like ChatGPT
			\item Generates plausible text without understanding
			\item Based on statistical patterns
		\end{itemize}
		
		\vspace{0.1cm}
		\textbf{CHARACTERISTICS:}
		\begin{itemize}
			\item Imitates responses from training data
			\item With variation
			\item No internal model of domain
			\item No understanding of what's discussed
		\end{itemize}
		
		\vspace{0.1cm}
		\textbf{EVIDENCE:}
		\begin{itemize}
			\item Chess example
			\item Common openings: sensible moves
			\item Uncommon openings: errors appear
			\item No board model
			\item No analysis or lookahead
		\end{itemize}
		
		\vspace{0.1cm}
		\textbf{IMPLICATIONS:}
		\begin{itemize}
			\item Apparent intelligence is illusion
			\item Cannot trust outputs
			\item Need for verification
			\item Hallucination inevitable
		\end{itemize}
		
		\textcolor{myred}{\textbf{EXAM TRAP:}} Stochastic parrot = \textcolor{myblue}{mimicking without understanding}!
	\end{frame}
	
	\begin{frame}{T74: Reinforcement Learning from Human Feedback}
		\fontsize{6.5}{7.5}\selectfont
		\textbf{TOPIC: Reinforcement Learning from Human Feedback (RLHF)}
		
		\vspace{0.1cm}
		\textbf{DEFINITION:}
		\begin{itemize}
			\item Training technique for LLMs
			\item Humans assess relative suitability of responses
			\item Feedback used to generate reward model
		\end{itemize}
		
		\vspace{0.1cm}
		\textbf{PROCESS:}
		\begin{enumerate}
			\item System generates range of responses for prompt
			\item Humans assess relative suitability
			\item Statistical analysis of feedback
			\item Generate "reward model"
			\item Use reward model to rate responses
			\item Select most appropriate responses
		\end{enumerate}
		
		\vspace{0.1cm}
		\textbf{APPLICATION:}
		\begin{itemize}
			\item ChatGPT training
			\item Improving response quality
			\item Aligning with human preferences
			\item Reducing harmful outputs
		\end{itemize}
		
		\vspace{0.1cm}
		\textbf{LIMITATIONS:}
		\begin{itemize}
			\item Based on human judgments
			\item Humans may have biases
			\item Cannot eliminate all issues
			\item Hallucination remains inevitable
		\end{itemize}
		
		\textcolor{myred}{\textbf{EXAM TRAP:}} RLHF = \textcolor{myblue}{using human feedback to improve LLMs}!
	\end{frame}
	
	\begin{frame}{T75: Supervised Fine-Tuning}
		\fontsize{6.5}{7.5}\selectfont
		\textbf{TOPIC: Supervised Fine-Tuning}
		
		\vspace{0.1cm}
		\textbf{DEFINITION:}
		\begin{itemize}
			\item Training technique for LLMs
			\item Human-crafted responses to typical prompts
			\item Supervised learning approach
		\end{itemize}
		
		\vspace{0.1cm}
		\textbf{PROCESS:}
		\begin{enumerate}
			\item Collect typical prompts
			\item Human experts craft suitable responses
			\item Train model on these examples
			\item Model learns to generate similar responses
		\end{enumerate}
		
		\vspace{0.1cm}
		\textbf{PURPOSE:}
		\begin{itemize}
			\item Improve response quality
			\item Align with human expectations
			\item Reduce inappropriate outputs
			\item Guide model behavior
		\end{itemize}
		
		\vspace{0.1cm}
		\textbf{RELATIONSHIP TO RLHF:}
		\begin{itemize}
			\item First stage of training
			\item Followed by RLHF
			\item Both use human input
			\item Complementary approaches
		\end{itemize}
		
		\textcolor{myred}{\textbf{EXAM TRAP:}} Fine-tuning = \textcolor{myblue}{learning from human-crafted examples}!
	\end{frame}
	
	\begin{frame}{T76: Principal Component Analysis - Definition}
		\fontsize{6.5}{7.5}\selectfont
		\textbf{TOPIC: Principal Component Analysis (PCA) - Definition}
		
		\vspace{0.1cm}
		\textbf{DEFINITION:}
		\begin{itemize}
			\item Classical method for dimensionality reduction
			\item Identifies main dimensions of variation
			\item Enables easier grasp of range and relationships
		\end{itemize}
		
		\vspace{0.1cm}
		\textbf{HOW IT WORKS:}
		\begin{itemize}
			\item Transform coordinates
			\item New dimensions align with best discrimination
			\item First vector: most discriminating
			\item Second vector: orthogonal, next most
			\item Continue for all dimensions
		\end{itemize}
		
		\vspace{0.1cm}
		\textbf{APPLICATIONS:}
		\begin{itemize}
			\item Village layout example
			\item Stylometry example
			\item Data visualization
			\item Feature extraction
		\end{itemize}
		
		\vspace{0.1cm}
		\textbf{ADVANTAGES:}
		\begin{itemize}
			\item Powerful dimensionality reduction
			\item Preserves significant information
			\item Highlights clusters
			\item Less inscrutable than deep learning
		\end{itemize}
		
		\textcolor{myred}{\textbf{EXAM TRAP:}} PCA = \textcolor{myblue}{finding most informative dimensions}!
	\end{frame}
	
	\begin{frame}{T77: Dimensionality Reduction}
		\fontsize{6.5}{7.5}\selectfont
		\textbf{TOPIC: Dimensionality Reduction}
		
		\vspace{0.1cm}
		\textbf{DEFINITION:}
		\begin{itemize}
			\item Technique in unsupervised learning
			\item Identifies main dimensions of variation
			\item Enables easier grasp of range and relationships
		\end{itemize}
		
		\vspace{0.1cm}
		\textbf{PURPOSE:}
		\begin{itemize}
			\item Simplify complex data
			\item Reduce computational requirements
			\item Visualize high-dimensional data
			\item Remove redundant features
		\end{itemize}
		
		\vspace{0.1cm}
		\textbf{METHODS:}
		\begin{itemize}
			\item Principal Component Analysis (PCA)
			\item Autoencoders
			\item t-SNE
			\item UMAP
		\end{itemize}
		
		\vspace{0.1cm}
		\textbf{PCA ADVANTAGE:}
		\begin{itemize}
			\item Less inscrutable than deep learning
			\item Understand what dimensions represent
			\item Explore basis straightforwardly
		\end{itemize}
		
		\textcolor{myred}{\textbf{EXAM TRAP:}} Dimensionality reduction = \textcolor{myblue}{simplifying complex data}!
	\end{frame}
	
	\begin{frame}{T78: Cluster Analysis}
		\fontsize{6.5}{7.5}\selectfont
		\textbf{TOPIC: Cluster Analysis}
		
		\vspace{0.1cm}
		\textbf{DEFINITION:}
		\begin{itemize}
			\item Learning technique in unsupervised learning
			\item Groups of similar objects identified
			\item Without training set
			\item Automatic identification
		\end{itemize}
		
		\vspace{0.1cm}
		\textbf{PURPOSE:}
		\begin{itemize}
			\item Discover natural groupings
			\item Segment data
			\item Identify patterns
			\item Exploratory analysis
		\end{itemize}
		
		\vspace{0.1cm}
		\textbf{METHODS:}
		\begin{itemize}
			\item K-means clustering
			\item Hierarchical clustering
			\item DBSCAN
			\item Gaussian mixture models
		\end{itemize}
		
		\vspace{0.1cm}
		\textbf{APPLICATIONS:}
		\begin{itemize}
			\item Customer segmentation
			\item Document clustering
			\item Image segmentation
			\item Anomaly detection
		\end{itemize}
		
		\textcolor{myred}{\textbf{EXAM TRAP:}} Cluster analysis = \textcolor{myblue}{automatic grouping of similar objects}!
	\end{frame}
	
	\begin{frame}{T79: Data Leakage}
		\fontsize{6.5}{7.5}\selectfont
		\textbf{TOPIC: Data Leakage}
		
		\vspace{0.1cm}
		\textbf{DEFINITION:}
		\begin{itemize}
			\item Using future or concurrent data
			\item Wouldn't have been observable at training time
			\item Causes optimistic performance estimates
		\end{itemize}
		
		\vspace{0.1cm}
		\textbf{EXAMPLE:}
		\begin{itemize}
			\item Predicting loan defaults
			\item Using data from after loan approval
			\item Model appears accurate
			\item But fails in production
		\end{itemize}
		
		\vspace{0.1cm}
		\textbf{PREVENTION:}
		\begin{itemize}
			\item Careful data splitting
			\item Time-based validation
			\item Cross-validation
			\item Domain knowledge
		\end{itemize}
		
		\vspace{0.1cm}
		\textbf{CONSEQUENCES:}
		\begin{itemize}
			\item Overly optimistic performance
			\item Poor generalization
			\item Failed deployments
			\item Financial losses
		\end{itemize}
		
		\textcolor{myred}{\textbf{EXAM TRAP:}} Data leakage = \textcolor{myblue}{using future data in training}!
	\end{frame}
	
	\begin{frame}{T80: Model Robustness}
		\fontsize{6.5}{7.5}\selectfont
		\textbf{TOPIC: Model Robustness}
		
		\vspace{0.1cm}
		\textbf{DEFINITION:}
		\begin{itemize}
			\item Measure of how well model performs
			\item With inputs previously not encountered
			\item Or changes in data
		\end{itemize}
		
		\vspace{0.1cm}
		\textbf{IMPORTANCE:}
		\begin{itemize}
			\item Real-world data varies
			\item Conditions change
			\item Edge cases exist
			\item Adversarial attacks
		\end{itemize}
		
		\vspace{0.1cm}
		\textbf{CHALLENGES:}
		\begin{itemize}
			\item Deep learning vulnerable to adversarial examples
			\item Small changes cause major errors
			\item Overconfidence in wrong predictions
			\item Difficulty detecting failures
		\end{itemize}
		
		\vspace{0.1cm}
		\textbf{IMPROVEMENT:}
		\begin{itemize}
			\item Extensive testing
			\item Adversarial training
			\item Ensemble methods
			\item Regularization
		\end{itemize}
		
		\textcolor{myred}{\textbf{EXAM TRAP:}} Robustness = \textcolor{myblue}{performing well on unseen data}!
	\end{frame}
	
	\begin{frame}{T81: Model Interpretability}
		\fontsize{6.5}{7.5}\selectfont
		\textbf{TOPIC: Model Interpretability}
		
		\vspace{0.1cm}
		\textbf{DEFINITION:}
		\begin{itemize}
			\item Degree to which human can comprehend model behavior
			\item With built-in mechanisms
			\item Understand how inputs affect outputs
		\end{itemize}
		
		\vspace{0.1cm}
		\textbf{SPECTRUM:}
		\begin{itemize}
			\item \textcolor{myblue}{High:} Decision trees, linear regression
			\item \textcolor{myblue}{Medium:} Random forests, gradient boosting
			\item \textcolor{myblue}{Low:} Deep neural networks
		\end{itemize}
		
		\vspace{0.1cm}
		\textbf{TRADE-OFF:}
		\begin{itemize}
			\item Complex models often more accurate
			\item But less interpretable
			\item Need to balance
			\item Depends on application
		\end{itemize}
		
		\vspace{0.1cm}
		\textbf{TOOLS:}
		\begin{itemize}
			\item Feature importance
			\item SHAP values
			\item LIME
			\item Counterfactual explanations
		\end{itemize}
		
		\textcolor{myred}{\textbf{EXAM TRAP:}} Interpretability vs. accuracy \textcolor{myblue}{trade-off}!
	\end{frame}
	
	\begin{frame}{T82: Model Explainability}
		\fontsize{6.5}{7.5}\selectfont
		\textbf{TOPIC: Model Explainability}
		
		\vspace{0.1cm}
		\textbf{DEFINITION:}
		\begin{itemize}
			\item Capacity to explain how AI makes decisions
			\item Often after the fact (ex-post)
			\item Understandable terms
		\end{itemize}
		
		\vspace{0.1cm}
		\textbf{IMPORTANCE:}
		\begin{itemize}
			\item Accountability
			\item Trust building
			\item Error identification
			\item Bias detection
		\end{itemize}
		
		\vspace{0.1cm}
		\textbf{CHALLENGES:}
		\begin{itemize}
			\item Deep networks opaque
			\item Billions of weights
			\item No simple explanation
			\item Hard to identify bias sources
		\end{itemize}
		
		\vspace{0.1cm}
		\textbf{APPROACHES:}
		\begin{itemize}
			\item Local explanations
			\item Global explanations
			\item Model-specific methods
			\item Model-agnostic methods
		\end{itemize}
		
		\textcolor{myred}{\textbf{EXAM TRAP:}} Explainability = \textcolor{myblue}{explaining decisions after the fact}!
	\end{frame}
	
	\begin{frame}{T83: Training Data Quality}
		\fontsize{6.5}{7.5}\selectfont
		\textbf{TOPIC: Training Data Quality}
		
		\vspace{0.1cm}
		\textbf{IMPORTANCE:}
		\begin{itemize}
			\item Models only as good as training data
			\item Garbage in, garbage out
			\item Quality determines performance
		\end{itemize}
		
		\vspace{0.1cm}
		\textbf{KEY ATTRIBUTES:}
		\begin{itemize}
			\item Accuracy
			\item Completeness
			\item Consistency
			\item Relevance
			\item Timeliness
		\end{itemize}
		
		\vspace{0.1cm}
		\textbf{COMMON ISSUES:}
		\begin{itemize}
			\item Missing values
			\item Outliers
			\item Noisy data
			\item Biased sampling
			\item Label errors
		\end{itemize}
		
		\vspace{0.1cm}
		\textbf{BEST PRACTICES:}
		\begin{itemize}
			\item Data validation
			\item Preprocessing
			\item Documentation
			\item Continuous monitoring
		\end{itemize}
		
		\textcolor{myred}{\textbf{EXAM TRAP:}} Training data quality \textcolor{myblue}{determines model quality}!
	\end{frame}
	
	\begin{frame}{T84: Feature Engineering}
		\fontsize{6.5}{7.5}\selectfont
		\textbf{TOPIC: Feature Engineering}
		
		\vspace{0.1cm}
		\textbf{DEFINITION:}
		\begin{itemize}
			\item Selecting and engineering appropriate features
			\item Variables used as inputs for model
			\item Critical for model performance
		\end{itemize}
		
		\vspace{0.1cm}
		\textbf{ACTIVITIES:}
		\begin{itemize}
			\item Transforming variables
			\item Combining variables
			\item Creating derived features
			\item Dimensionality reduction
		\end{itemize}
		
		\vspace{0.1cm}
		\textbf{EXAMPLES:}
		\begin{itemize}
			\item Log transformations
			\item Polynomial features
			\item Interaction terms
			\item Domain-specific ratios
		\end{itemize}
		
		\vspace{0.1cm}
		\textbf{IMPORTANCE:}
		\begin{itemize}
			\item Good features improve performance
			\item Reduces need for complex models
			\item Domain knowledge valuable
			\item Iterative process
		\end{itemize}
		
		\textcolor{myred}{\textbf{EXAM TRAP:}} Feature engineering = \textcolor{myblue}{creating informative inputs}!
	\end{frame}
	
	\begin{frame}{T85: Model Selection}
		\fontsize{6.5}{7.5}\selectfont
		\textbf{TOPIC: Model Selection}
		
		\vspace{0.1cm}
		\textbf{DEFINITION:}
		\begin{itemize}
			\item Choosing modeling approach
			\item Comprehensive evaluation of alternatives
			\item Documenting rationale
		\end{itemize}
		
		\vspace{0.1cm}
		\textbf{CONSIDERATIONS:}
		\begin{itemize}
			\item Problem type (classification, regression)
			\item Data characteristics
			\item Interpretability requirements
			\item Computational resources
			\item Regulatory requirements
		\end{itemize}
		
		\vspace{0.1cm}
		\textbf{APPROACHES:}
		\begin{itemize}
			\item Start simple
			\item Compare alternatives
			\item Cross-validation
			\item Consider trade-offs
		\end{itemize}
		
		\vspace{0.1cm}
		\textbf{DOCUMENTATION:}
		\begin{itemize}
			\item Rationale for selection
			\item Alternatives considered
			\item Performance comparison
			\item Limitations
		\end{itemize}
		
		\textcolor{myred}{\textbf{EXAM TRAP:}} Model selection requires \textcolor{myblue}{systematic evaluation}!
	\end{frame}
	
	\begin{frame}{T86: Hyperparameter Tuning}
		\fontsize{6.5}{7.5}\selectfont
		\textbf{TOPIC: Hyperparameter Tuning}
		
		\vspace{0.1cm}
		\textbf{DEFINITION:}
		\begin{itemize}
			\item Optimizing parameters not learned from data
			\item Set before training
			\item Configuration settings
		\end{itemize}
		
		\vspace{0.1cm}
		\textbf{EXAMPLES:}
		\begin{itemize}
			\item Learning rate
			\item Number of hidden layers
			\item Regularization parameter
			\item Number of epochs
		\end{itemize}
		
		\vspace{0.1cm}
		\textbf{METHODS:}
		\begin{itemize}
			\item Grid search
			\item Random search
			\item Bayesian optimization
			\item Cross-validation
		\end{itemize}
		
		\vspace{0.1cm}
		\textbf{IMPORTANCE:}
		\begin{itemize}
			\item Affects model performance
			\item Balance bias-variance
			\item Avoid overfitting
			\item Optimize generalization
		\end{itemize}
		
		\textcolor{myred}{\textbf{EXAM TRAP:}} Hyperparameters = \textcolor{myblue}{settings chosen before training}!
	\end{frame}
	
	\begin{frame}{T87: Cross-Validation}
		\fontsize{6.5}{7.5}\selectfont
		\textbf{TOPIC: Cross-Validation}
		
		\vspace{0.1cm}
		\textbf{DEFINITION:}
		\begin{itemize}
			\item Technique for model evaluation
			\item Split data into folds
			\item Train on some, validate on others
			\item Repeat and average
		\end{itemize}
		
		\vspace{0.1cm}
		\textbf{K-FOLD CROSS-VALIDATION:}
		\begin{itemize}
			\item Split data into k folds
			\item Train on k-1 folds
			\item Validate on 1 fold
			\item Repeat k times
			\item Average results
		\end{itemize}
		
		\vspace{0.1cm}
		\textbf{BENEFITS:}
		\begin{itemize}
			\item Efficient use of data
			\item Reliable performance estimate
			\item Reduces overfitting
			\item Model selection
		\end{itemize}
		
		\vspace{0.1cm}
		\textbf{VARIANTS:}
		\begin{itemize}
			\item Stratified (preserves class distribution)
			\item Leave-one-out (LOOCV)
			\item Time series (preserves order)
		\end{itemize}
		
		\textcolor{myred}{\textbf{EXAM TRAP:}} Cross-validation = \textcolor{myblue}{train on k-1, validate on 1}!
	\end{frame}
	
	\begin{frame}{T88: Model Testing Types}
		\fontsize{6.5}{7.5}\selectfont
		\textbf{TOPIC: Model Testing Types}
		
		\vspace{0.1cm}
		\textbf{CHRONOLOGICAL ORDER:}
		\begin{enumerate}
			\item \textcolor{myblue}{Unit Tests:} Individual code pieces
			\item \textcolor{myblue}{Component Tests:} Individual code sections
			\item \textcolor{myblue}{Integration Tests:} Interfaces between components
			\item \textcolor{myblue}{System Tests:} Complete integrated system
			\item \textcolor{myblue}{Performance Tests:} Speed and stability
			\item \textcolor{myblue}{Regression Tests:} After modifications
			\item \textcolor{myblue}{Acceptance Tests:} Ready for deployment
			\item \textcolor{myblue}{Adversarial Tests:} Extreme scenarios
		\end{enumerate}
		
		\vspace{0.1cm}
		\textbf{TESTING APPROACHES:}
		\begin{itemize}
			\item White-box: access to code
			\item Black-box: user perspective
			\item Gray-box: knowledge of architecture
		\end{itemize}
		
		\vspace{0.1cm}
		\textbf{KEY PRINCIPLES:}
		\begin{itemize}
			\item Start testing early
			\item Use realistic configurations
			\item Less realistic parameters better
			\item Continuous testing
		\end{itemize}
		
		\textcolor{myred}{\textbf{EXAM TRAP:}} Testing should be \textcolor{myblue}{comprehensive and start early}!
	\end{frame}
	
	\begin{frame}{T89: White-Box Testing}
		\fontsize{6.5}{7.5}\selectfont
		\textbf{TOPIC: White-Box Testing}
		
		\vspace{0.1cm}
		\textbf{DEFINITION:}
		\begin{itemize}
			\item Testers have access to internal data structures
			\item Access to algorithms
			\item Access to actual code
		\end{itemize}
		
		\vspace{0.1cm}
		\textbf{CHARACTERISTICS:}
		\begin{itemize}
			\item Testing at unit level
			\item From developer's perspective
			\item Line-by-line proofreading
			\item More effective for finding specific issues
		\end{itemize}
		
		\vspace{0.1cm}
		\textbf{WHEN TO USE:}
		\begin{itemize}
			\item Internally developed models
			\item Early stage of development
			\item When code access available
		\end{itemize}
		
		\vspace{0.1cm}
		\textbf{ADVANTAGES:}
		\begin{itemize}
			\item Thorough testing
			\item Can identify specific issues
			\item Deeper understanding
		\end{itemize}
		
		\vspace{0.1cm}
		\textbf{DISADVANTAGES:}
		\begin{itemize}
			\item May introduce bias
			\item Limited to accessible code
			\item Not suitable for proprietary models
		\end{itemize}
		
		\textcolor{myred}{\textbf{EXAM TRAP:}} White-box = \textcolor{myblue}{access to internal code}!
	\end{frame}
	
	\begin{frame}{T90: Black-Box Testing}
		\fontsize{6.5}{7.5}\selectfont
		\textbf{TOPIC: Black-Box Testing}
		
		\vspace{0.1cm}
		\textbf{DEFINITION:}
		\begin{itemize}
			\item Treats software as closed box
			\item From user's perspective
			\item No knowledge of internal implementation
		\end{itemize}
		
		\vspace{0.1cm}
		\textbf{CHARACTERISTICS:}
		\begin{itemize}
			\item Reduces likelihood of bias
			\item Functional testing
			\item Focus on inputs and outputs
			\item User experience testing
		\end{itemize}
		
		\vspace{0.1cm}
		\textbf{WHEN TO USE:}
		\begin{itemize}
			\item Proprietary vendor models
			\item Functional testing
			\item Final stages of development
			\item When code not accessible
		\end{itemize}
		
		\vspace{0.1cm}
		\textbf{ADVANTAGES:}
		\begin{itemize}
			\item Reduces bias
			\item Tests actual user experience
			\item Suitable for all models
		\end{itemize}
		
		\vspace{0.1cm}
		\textbf{DISADVANTAGES:}
		\begin{itemize}
			\item Less thorough
			\item Harder to identify specific issues
			\item Cannot see internal workings
		\end{itemize}
		
		\textcolor{myred}{\textbf{EXAM TRAP:}} Black-box = \textcolor{myblue}{user perspective, no code access}!
	\end{frame}
	
	\begin{frame}{T91: Gray-Box Testing}
		\fontsize{6.5}{7.5}\selectfont
		\textbf{TOPIC: Gray-Box Testing}
		
		\vspace{0.1cm}
		\textbf{DEFINITION:}
		\begin{itemize}
			\item Between white-box and black-box
			\item Informed by knowledge of functional specifications
			\item Knowledge of architecture
			\item Does not access code itself
		\end{itemize}
		
		\vspace{0.1cm}
		\textbf{CHARACTERISTICS:}
		\begin{itemize}
			\item Combines benefits of both approaches
			\item Knowledge of system design
			\item No direct code access
			\item Balanced perspective
		\end{itemize}
		
		\vspace{0.1cm}
		\textbf{WHEN TO USE:}
		\begin{itemize}
			\item Distributed applications
			\item Testing component integration
			\item When architecture known
			\item Complex systems
		\end{itemize}
		
		\vspace{0.1cm}
		\textbf{ADVANTAGES:}
		\begin{itemize}
			\item More informed than black-box
			\item Less biased than white-box
			\item Flexible approach
		\end{itemize}
		
		\textcolor{myred}{\textbf{EXAM TRAP:}} Gray-box = \textcolor{myblue}{knowledge of architecture, not code}!
	\end{frame}
	
	\begin{frame}{T92: Use Test}
		\fontsize{6.5}{7.5}\selectfont
		\textbf{TOPIC: The Use Test}
		
		\vspace{0.1cm}
		\textbf{DEFINITION:}
		\begin{itemize}
			\item Examines model within context
			\item Considers human interaction
			\item Assesses actual usage
			\item Evaluates acceptance
		\end{itemize}
		
		\vspace{0.1cm}
		\textbf{KEY QUESTIONS:}
		\begin{itemize}
			\item Are results utilized appropriately?
			\item Do they prove useful?
			\item If not, why?
		\end{itemize}
		
		\vspace{0.1cm}
		\textbf{FOCUS AREAS:}
		\begin{itemize}
			\item People, not numbers
			\item Acceptance and trust
			\item Comprehensibility
			\item Communication
		\end{itemize}
		
		\vspace{0.1cm}
		\textbf{REPORTING:}
		\begin{itemize}
			\item Presented to senior management
			\item Not technical reports
			\item Qualitative in nature
			\item Improves culture
		\end{itemize}
		
		\textcolor{myred}{\textbf{EXAM TRAP:}} Use test focuses on \textcolor{myblue}{how people use the model}!
	\end{frame}
	
	\begin{frame}{T93: Model Limitations}
		\fontsize{6.5}{7.5}\selectfont
		\textbf{TOPIC: Model Limitations}
		
		\vspace{0.1cm}
		\textbf{DEFINITION:}
		\begin{itemize}
			\item Constraints on model applicability
			\item Boundaries beyond which model invalid
			\item Should be documented
			\item Must be communicated to users
		\end{itemize}
		
		\vspace{0.1cm}
		\textbf{TYPES:}
		\begin{itemize}
			\item \textcolor{myblue}{Product Applicability:} Specific products only
			\item \textcolor{myblue}{Access Restrictions:} Who can use
			\item \textcolor{myblue}{Portfolio Size Limits:} For new markets
			\item \textcolor{myblue}{Data Limitations:} Quality, coverage
		\end{itemize}
		
		\vspace{0.1cm}
		\textbf{WHY LIMITS:}
		\begin{itemize}
			\item Prevent misuse
			\item Ensure appropriate application
			\item Maintain risk ranking
			\item Manage expectations
		\end{itemize}
		
		\vspace{0.1cm}
		\textbf{DOCUMENTATION:}
		\begin{itemize}
			\item In model inventory
			\item In model documentation
			\item Communicated to users
			\item Regular review
		\end{itemize}
		
		\textcolor{myred}{\textbf{EXAM TRAP:}} Model limits prevent \textcolor{myblue}{misuse beyond intended application}!
	\end{frame}
	
	\begin{frame}{T94: Model Risk Ranking}
		\fontsize{6.5}{7.5}\selectfont
		\textbf{TOPIC: Model Risk Ranking}
		
		\vspace{0.1cm}
		\textbf{DEFINITION:}
		\begin{itemize}
			\item Classification of model risk
			\item Drives validation depth and frequency
			\item Based on multiple factors
		\end{itemize}
		
		\vspace{0.1cm}
		\textbf{KEY FACTORS:}
		\begin{enumerate}
			\item \textcolor{myblue}{Materiality:}
			\begin{itemize}
				\item Economic consequences
				\item Operational impact
				\item Reputational consequences
			\end{itemize}
			\item \textcolor{myblue}{Complexity:}
			\begin{itemize}
				\item May elevate ranking
				\item Even if materiality low
			\end{itemize}
			\item \textcolor{myblue}{Exposure:}
			\begin{itemize}
				\item Published financial reports
				\item Wide impact
			\end{itemize}
		\end{enumerate}
		
		\vspace{0.1cm}
		\textbf{IMPORTANT:}
		\begin{itemize}
			\item Minimum standards for all models
			\item Regardless of materiality
			\item Absence of assessment is risk
		\end{itemize}
		
		\textcolor{myred}{\textbf{EXAM TRAP:}} Risk ranking considers \textcolor{myblue}{materiality, complexity, exposure}!
	\end{frame}
	
	\begin{frame}{T95: Three Lines of Defense}
		\fontsize{6.5}{7.5}\selectfont
		\textbf{TOPIC: Three Lines of Defense}
		
		\vspace{0.1cm}
		\textbf{FIRST LINE:}
		\begin{itemize}
			\item Business and model owners/developers
			\item Own the model and performance
			\item Design and implement models
			\item First to react to breaches
		\end{itemize}
		
		\vspace{0.1cm}
		\textbf{SECOND LINE:}
		\begin{itemize}
			\item Independent risk management/validation
			\item Set policy
			\item Conduct validations
			\item Challenge First Line
		\end{itemize}
		
		\vspace{0.1cm}
		\textbf{THIRD LINE:}
		\begin{itemize}
			\item Internal audit
			\item Independent assurance
			\item Review framework effectiveness
			\item Ensure compliance
		\end{itemize}
		
		\vspace{0.1cm}
		\textbf{COMMUNICATION:}
		\begin{itemize}
			\item Breach notification
			\item Independent review
			\item Dispute resolution
			\item Reporting
		\end{itemize}
		
		\textcolor{myred}{\textbf{EXAM TRAP:}} Three lines provide \textcolor{myblue}{checks and balances}!
	\end{frame}
	
	\begin{frame}{T96: Model Validation}
		\fontsize{6.5}{7.5}\selectfont
		\textbf{TOPIC: Model Validation}
		
		\vspace{0.1cm}
		\textbf{DEFINITION:}
		\begin{itemize}
			\item Rigorously testing model, inputs, outputs, governance
			\item Identify extent of residual model risk
			\item Assess mitigating controls
		\end{itemize}
		
		\vspace{0.1cm}
		\textbf{WHEN PERFORMED:}
		\begin{itemize}
			\item \textcolor{myblue}{Initial:} Before implementation
			\item \textcolor{myblue}{Periodic:} Annual or more frequent
			\item \textcolor{myblue}{Change-Based:} After material changes
		\end{itemize}
		
		\vspace{0.1cm}
		\textbf{KEY ACTIVITIES:}
		\begin{itemize}
			\item Data assessment
			\item Assumptions review
			\item Limitations identification
			\item Implementation review
			\item Performance monitoring
			\item Documentation review
		\end{itemize}
		
		\vspace{0.1cm}
		\textbf{KEY PRINCIPLE:}
		\begin{itemize}
			\item No "valid" model
			\item Attempt to invalidate
			\item Usefulness vs. validity
			\item Effective challenge
		\end{itemize}
		
		\textcolor{myred}{\textbf{EXAM TRAP:}} Validation = \textcolor{myblue}{attempt to invalidate, not prove validity}!
	\end{frame}
	
	\begin{frame}{T97: Model Changes}
		\fontsize{6.5}{7.5}\selectfont
		\textbf{TOPIC: Model Changes}
		
		\vspace{0.1cm}
		\textbf{CHANGE MANAGEMENT:}
		\begin{itemize}
			\item Structured process
			\item Change requests
			\item Impact assessment
			\item Documentation
			\item Validation
			\item Approval
		\end{itemize}
		
		\vspace{0.1cm}
		\textbf{INDEPENDENT REVIEW:}
		\begin{itemize}
			\item By individuals not involved
			\item Assess accuracy, reliability
			\item Compliance
			\item Potential risks
		\end{itemize}
		
		\vspace{0.1cm}
		\textbf{AI/ML DIFFERENCES:}
		\begin{itemize}
			\item More frequent monitoring
			\item Broader stakeholders
			\item Focus on explainability
			\item Benchmarking
			\item Different data sets
			\item Specific cases
		\end{itemize}
		
		\vspace{0.1cm}
		\textbf{IMPLEMENTATION:}
		\begin{itemize}
			\item Controlled manner
			\item Documentation
			\item Monitoring changes
			\item Automated controls
		\end{itemize}
		
		\textcolor{myred}{\textbf{EXAM TRAP:}} Model changes require \textcolor{myblue}{structured governance}!
	\end{frame}
	
	\begin{frame}{T98: Data Governance}
		\fontsize{6.5}{7.5}\selectfont
		\textbf{TOPIC: Data Governance}
		
		\vspace{0.1cm}
		\textbf{DEFINITION:}
		\begin{itemize}
			\item System of decision rights and accountabilities
			\item For information-related processes
			\item People, processes, technologies
		\end{itemize}
		
		\vspace{0.1cm}
		\textbf{KEY COMPONENTS:}
		\begin{enumerate}
			\item \textcolor{myblue}{Strategy:} Vision, goals, policies
			\item \textcolor{myblue}{Quality:} Accuracy, consistency, integrity
			\item \textcolor{myblue}{Provenance:} Legal right to use
			\item \textcolor{myblue}{Classification:} Structure, confidentiality
			\item \textcolor{myblue}{Metadata:} Documentation, dictionary
			\item \textcolor{myblue}{Security:} Protection, compliance
		\end{enumerate}
		
		\vspace{0.1cm}
		\textbf{THREE PRINCIPLES:}
		\begin{itemize}
			\item Obtaining consent
			\item Minimizing data held
			\item Ensuring rights of data subjects
		\end{itemize}
		
		\vspace{0.1cm}
		\textbf{BOARD ROLE:}
		\begin{itemize}
			\item Approve framework
			\item Oversee compliance
			\item Assess risks
			\item Regular review
		\end{itemize}
		
		\textcolor{myred}{\textbf{EXAM TRAP:}} Data governance = \textcolor{myblue}{decision rights and accountabilities}!
	\end{frame}
	
	\begin{frame}{T99: GenAI Governance}
		\fontsize{6.5}{7.5}\selectfont
		\textbf{TOPIC: GenAI Governance}
		
		\vspace{0.1cm}
		\textbf{WHY HARDER:}
		\begin{itemize}
			\item Opacity and unpredictability
			\item Emergent capabilities
			\item Post-deployment shifts
			\item Embedded in workflows
			\item Fragmented regulation
		\end{itemize}
		
		\vspace{0.1cm}
		\textbf{INTERNAL LEVERS:}
		\begin{itemize}
			\item Assign ownership
			\item Risk-based review
			\item Incentivize responsible behavior
			\item Staff training
			\item Escalation paths
		\end{itemize}
		
		\vspace{0.1cm}
		\textbf{EXTERNAL CONSTRAINTS:}
		\begin{itemize}
			\item Track regulations
			\item Evaluate provider responsibilities
			\item Align with standards
		\end{itemize}
		
		\vspace{0.1cm}
		\textbf{ADAPTIVE TOOLS:}
		\begin{itemize}
			\item Continuous monitoring
			\item Red-teaming
			\item Scenario planning
			\item Technical safeguards
			\item Feedback loops
		\end{itemize}
		
		\textcolor{myred}{\textbf{EXAM TRAP:}} GenAI governance requires \textcolor{myblue}{internal, external, adaptive approaches}!
	\end{frame}
	
	\begin{frame}{T100: AI and Risk - Final Summary}
		\fontsize{6.5}{7.5}\selectfont
		\textbf{TOPIC: AI and Risk - Comprehensive Summary}
		
		\vspace{0.1cm}
		\textbf{CLASSICAL AI:}
		\begin{itemize}
			\item Explicit, rule-based processing
			\item Search, recursion, lookahead
			\item Limits: brittleness, frame problem
		\end{itemize}
		
		\vspace{0.1cm}
		\textbf{NEURAL NETWORKS:}
		\begin{itemize}
			\item Biologically inspired
			\item Deep learning, backpropagation
			\item Inscrutability challenges
		\end{itemize}
		
		\vspace{0.1cm}
		\textbf{MACHINE LEARNING TYPES:}
		\begin{itemize}
			\item Reinforcement, supervised, unsupervised, semi-supervised
		\end{itemize}
		
		\vspace{0.1cm}
		\textbf{RISKS:}
		\begin{itemize}
			\item Inscrutability: hidden bias, privacy
			\item Over-reliance: automation bias
			\item Individual, organizational, societal
		\end{itemize}
		
		\vspace{0.1cm}
		\textbf{KEY TAKEAWAYS:}
		\begin{itemize}
			\item AI is powerful but risky
			\item Understanding limitations crucial
			\item Responsible governance essential
			\item Balance innovation with risk management
		\end{itemize}
		
		\textcolor{myred}{\textbf{EXAM SUCCESS:}} Understand \textcolor{myblue}{classical AI, neural networks, ML types, and AI risks}!
	\end{frame}
	
\end{document}